\documentclass{article}

% PACKAGES

% correct encoding and typography
\usepackage[english]{babel}
\usepackage[utf8]{inputenc}
\usepackage[T1]{fontenc}

% math
\usepackage{amsfonts, latexsym, mathtools, amsthm, amssymb, amscd}
\usepackage{euscript}
\usepackage{esvect}
\usepackage{bm}
\usepackage{physics}

% aesthetics color options
\usepackage{xcolor, graphicx}
\usepackage{pdfpages}

% to access the last page
\usepackage{lastpage}

% better document formatting
\usepackage{fancyhdr}
\usepackage[a4paper, right=1in, left=1in, bottom=1.2in, top=1.2in, centering]{geometry}
\usepackage{fullpage}

\usepackage{datetime2}

% better stylistic formatting
\usepackage{parskip}
\usepackage{enumitem}
\usepackage{framed}
\usepackage{longtable}

% embedding and referencing
\usepackage{hyperref}
\hypersetup{colorlinks=true,citecolor=blue,urlcolor =black,linkbordercolor={1 0 0}}

% TikZ (workflow diagrams)
\usepackage{tikz}
\usetikzlibrary{arrows.meta,positioning,fit,calc,backgrounds,matrix}

%%%%%%%%%%%%%%%%%%%%%%%%%%%%%%%%%%%%%%%%%%%%%%%%%%%%%%%%%%%%%%%%%%%%%%%%%%%%%%%%%%%%%%%%%%%%%%%%%%%%%%%%%%%%%%%%%%
% Command formatting for ease-of-life

% new math symbols taking no arguments
\newcommand{\kk}{\Bbbk}
\newcommand{\minus}{\smallsetminus}
\newcommand{\dotp}{\boldsymbol{\cdot}}

%redefined math symbols taking no arguments
\newcommand{\<}{\langle}
\renewcommand{\>}{\rangle}
\renewcommand{\iff}{\Leftrightarrow}
\renewcommand{\implies}{\Rightarrow}

% Arrows
\newcommand{\from}{\leftarrow}
\newcommand{\ffrom}{\longleftarrow}
\newcommand{\goesto}{\rightsquigarrow}
\newcommand{\into}{\hookrightarrow}

% Basic bbface (blackboard bold) lettering command:
\newcommand{\bbb}[1]{\ensuremath{\mathbb{#1}}}

% Boldface lettering command used for vectors:
\newcommand{\bvec}[1]{\ensuremath{\mathbf{#1}}}

% Other lettering
\newcommand{\eps}{\varepsilon}


% some operators and delimiters
\DeclareMathOperator{\proj}{proj}
\DeclareMathOperator{\im}{im}
\DeclareMathOperator{\coker}{coker}
\DeclareMathOperator{\coim}{coim}
\DeclareMathOperator{\Span}{Span}
\DeclareMathOperator{\Stab}{Stab}
\DeclareMathOperator{\Orb}{Orb}
\DeclareMathOperator{\Cov}{Cov}
\DeclareMathOperator{\Var}{Var}
\DeclareMathOperator{\EE}{E}
\DeclareMathOperator{\Tor}{Tor}
\DeclareMathOperator{\Ext}{Ext}


% Other more involved shortcuts

%for overline math
\newcommand{\ol}[1]{{\overline{#1}}}

%redefined math symbols taking arguments
\renewcommand{\mod}[1]{\ (\mathrm{mod}\ #1)}

%for easy 2 x 2 matrices
\newcommand{\twobytwo}[1]{\left[\begin{array}{@{}cc@{}}#1\end{array}\right]}

%for easy column vectors of size 2
\newcommand{\tworow}[1]{\left[\begin{array}{@{}c@{}}#1\end{array}\right]}

%%%%%%%%%%%%%%%%%%%%%%%%%%%%%%%%%%%%%%%%%%%%%%%%%%%%%%%%%%%%%%%%%%%%%%%%%%%%%%%%%%%%%%%%%%%%%%%%%%%%%%%%%%%%%%%%%%
%Below are the theorem, definition, example, lemma, etc. body types.

\newtheorem{theorem}{Theorem}
\newtheorem*{proposition}{Proposition}
\newtheorem{lemma}[theorem]{Lemma}
\newtheorem{corollary}[theorem]{Corollary}
\newtheorem{conjecture}[theorem]{Conjecture}
\newtheorem{postulate}[theorem]{Postulate}
\theoremstyle{definition}
\newtheorem{defn}[theorem]{Definition}
\newtheorem{example}[theorem]{Example}

\theoremstyle{remark}
\newtheorem*{remark}{Remark}
\newtheorem*{notation}{Notation}
\newtheorem*{note}{Note}

%%%%%%%%%%%%%%%%%%%%%%%%%%%%%%%%%%%%%%%%%%%%%%%%%%%%%%%%%%%%%%%%%%%%%%%%%%%%%%%%%%%%%%%%%%%%%%%%%%%%%%%%%%%%%%%%%%
%Some formatting tidbits for margins, paragraphs, and removing orphans

% make widows and orphans rare
\clubpenalty =10000
\widowpenalty =10000

% formatting of paragraphs and separation
\setlength{\parindent}{0pt}
\setlength{\parskip}{5pt plus 1pt}
\setlength{\headheight}{13.6pt}

\setlength\marginparwidth{2.2in}
\setlength\marginparsep{1mm}

%%%%%%%%%%%%%%%%%%%%%%%%%%%%%%%%%%%%%%%%%%%%%%%%%%%%%%%%%

\title{Implementation and Ramifications of TAMER-OP: a Toolkit for Algebraic Module Encodings over $\mathbb{R}^n$ and Other Posets}
\author{Erik Mendes Novak}

\begin{document}
	\maketitle
	
	\tableofcontents
	
	
	\section{Introduction}
	
	%%%%%%%%%%%%%%%%%%%%%%%%%%%%%%%%%%%%%%%%%%%%%%%%%%%%%%%%%%%%%%%%%%%%%%%%%%%%%%%%
	% Section 2: Background mathematical information
	%%%%%%%%%%%%%%%%%%%%%%%%%%%%%%%%%%%%%%%%%%%%%%%%%%%%%%%%%%%%%%%%%%%%%%%%%%%%%%%%
	
	\section{Background mathematical information}\label{sec:background}
	
	Persistent homology assigns to data a family of homology groups indexed by a
	(possibly multi-dimensional) parameter, together with linear maps relating the
	groups as the parameter increases.  The resulting algebraic object is a
	\emph{persistence module}.  In one parameter, persistence modules admit a clean
	structure theory (barcodes) and robust metrics; in multiple parameters, the same
	classification program becomes substantially more subtle.  This section records
	the minimal language needed to state the objects studied in this thesis and to
	motivate the computational strategy pursued by TAMER-OP.
	
	Throughout, fix a coefficient field $\kk$ unless explicitly stated otherwise.
	When we discuss implementation choices later, we will emphasize why exact
	arithmetic over $\bbb{Q}$ is often preferred for certifying ranks and homological
	invariants (cf.\ \S\ref{sec:background:conventions}).
	
	\subsection{Posets as categories; persistence modules as functors}
	\label{sec:background:posets-as-cats}
	
	\subsubsection*{Posets and thin categories}
	
	\begin{defn}[Poset]\label{def:poset}
		A \emph{partially ordered set} (poset) is a set $Q$ equipped with a binary relation
		$\leq$ that is reflexive ($q\leq q$), antisymmetric ($q\leq q'$ and $q'\leq q$
		implies $q=q'$), and transitive ($q\leq q'$ and $q'\leq q''$ implies $q\leq q''$).
	\end{defn}
	
	\begin{defn}[Poset as a category]\label{def:poset-as-category}
		Given a poset $(Q,\leq)$, define a (small) category, which we also denote by $Q$,
		whose objects are elements of $Q$, and such that for $q,q'\in Q$ there is a unique
		morphism $q\to q'$ if and only if $q\leq q'$, and no morphism otherwise.  Composition
		is forced by transitivity, and identities are forced by reflexivity.
	\end{defn}
	
	A category arising from a poset in this way is \emph{thin}: between any two objects
	there is at most one morphism.  This viewpoint is the standard bridge between order
	theory and functorial algebraic constructions.
	
	\subsubsection*{Modules over a poset}
	
	\begin{defn}[Poset module / persistence module]\label{def:Q-module}
		Let $(Q,\leq)$ be a poset and let $\mathbf{Vect}_{\kk}$ denote the category of
		$\kk$-vector spaces and $\kk$-linear maps.  A \emph{$Q$-module} (or \emph{persistence
			module indexed by $Q$}) is a covariant functor
		\[
		M \colon Q \longrightarrow \mathbf{Vect}_{\kk}.
		\]
		Concretely, this is the data of:
		\begin{enumerate}[label=(\roman*)]
			\item a vector space $M(q)$ for each $q\in Q$, and
			\item for each comparable pair $q\leq q'$, a linear \emph{structure map}
			\[
			M(q\leq q') \colon M(q)\to M(q'),
			\]
		\end{enumerate}
		subject to the functoriality conditions
		\[
		M(q\leq q)=\mathrm{id}_{M(q)}\qquad\text{and}\qquad
		M(q\leq q'')=M(q'\leq q'')\circ M(q\leq q')\ \text{for }q\leq q'\leq q''.
		\]
	\end{defn}
	
	The indexing poset $Q$ will typically be $\bbb{R}$ or $\bbb{R}^n$ with the usual
	order(s), or $\bbb{Z}^n$ with the product (coordinatewise) order.  The degree $q$
	is sometimes called a \emph{parameter value}, and $M(q)$ a \emph{fiber} or
	\emph{stalk}.
	
	\begin{defn}[Morphisms]\label{def:morphism-of-Q-modules}
		A morphism of $Q$-modules $f\colon M\to N$ is a natural transformation of functors.
		Equivalently, it is a family of linear maps $f_q\colon M(q)\to N(q)$ indexed by
		$q\in Q$ such that for every $q\leq q'$ the following square commutes:
		\[
		\begin{CD}
			M(q) @>{M(q\leq q')}>> M(q')\\
			@V{f_q}VV @VV{f_{q'}}V\\
			N(q) @>{N(q\leq q')}>> N(q').
		\end{CD}
		\]
	\end{defn}
	
	We write $\mathrm{Hom}_Q(M,N)$ for the $\kk$-vector space of such natural
	transformations (when the fibers are finite-dimensional, $\mathrm{Hom}_Q(M,N)$
	itself is often computable by linear algebra).
	
	\subsubsection*{Functor categories are abelian}
	
	The category of $\kk$-vector spaces is abelian, and functor categories into an
	abelian category remain abelian.  In particular, $Q$-modules admit kernels,
	cokernels, images, and short exact sequences.
	
	\begin{proposition}[Pointwise abelian structure]\label{prop:functor-category-abelian}
		For any small category $\mathcal{C}$, the functor category
		$\mathbf{Vect}_{\kk}^{\mathcal{C}}=\mathrm{Fun}(\mathcal{C},\mathbf{Vect}_{\kk})$
		is abelian.  In particular, for $\mathcal{C}=Q$ a poset category, kernels and
		cokernels are computed \emph{objectwise}: for a morphism $f\colon M\to N$ one has
		\[
		(\ker f)(q)=\ker(f_q),\qquad (\mathrm{coker}\,f)(q)=\mathrm{coker}(f_q),
		\]
		and similarly for images and coimages, with the induced structure maps coming from
		functoriality.
	\end{proposition}
	
	\begin{remark}
		Proposition~\ref{prop:functor-category-abelian} is conceptually important because it
		reduces many categorical constructions on persistence modules to familiar linear
		algebra at each parameter value.  Much of TAMER-OP is built around exploiting this
		pointwise nature after reducing an a priori infinite indexing poset to a finite one.
	\end{remark}
	
	\subsubsection*{Upsets, downsets, and indicator modules}
	
	A recurrent theme (central to Miller's theory of tameness) is that many
	computationally tractable persistence modules can be assembled from very simple
	``indicator'' modules associated to order ideals and filters.
	
	\begin{defn}[Upsets and downsets]\label{def:upset-downset}
		Let $Q$ be a poset.
		\begin{enumerate}[label=(\roman*)]
			\item An \emph{upset} (order filter) is a subset $U\subseteq Q$ such that if
			$q\in U$ and $q\leq q'$, then $q'\in U$.
			\item A \emph{downset} (order ideal) is a subset $D\subseteq Q$ such that if
			$q'\in D$ and $q\leq q'$, then $q\in D$.
		\end{enumerate}
		For $q\in Q$, the \emph{principal upset} and \emph{principal downset} are
		\[
		\uparrow q := \{q'\in Q : q\leq q'\},\qquad
		\downarrow q := \{q'\in Q : q'\leq q\}.
		\]
	\end{defn}
	
	\begin{defn}[Indicator modules]\label{def:indicator-modules}
		Let $U\subseteq Q$ be an upset.  The \emph{upset indicator module} $\kk\{U\}$ is the
		$Q$-module defined by
		\[
		\kk\{U\}(q) :=
		\begin{cases}
			\kk & q\in U,\\
			0 & q\notin U,
		\end{cases}
		\]
		and for $q\leq q'$ the structure map $\kk\{U\}(q\leq q')$ is the identity
		$\kk\to\kk$ when $q,q'\in U$ and is the unique zero map otherwise.  Likewise, for a
		downset $D\subseteq Q$, define the \emph{downset indicator module} $\kk\{D\}$ by the
		same formula with $D$ in place of $U$.
	\end{defn}
	
	\begin{remark}
		The closure property in Definition~\ref{def:upset-downset} is exactly what guarantees
		functoriality for the ``identity within the region, zero outside'' convention in
		Definition~\ref{def:indicator-modules}.  For general subsets $S\subseteq Q$, the same
		prescription can fail to define a functor because composition along chains may force
		a map to be zero even when both endpoints lie in $S$.
	\end{remark}
	
	\begin{remark}[Yoneda intuition]\label{rmk:yoneda-principal-upset}
		When $Q$ is viewed as a category, the representable functor
		$\mathrm{Hom}_Q(q,-)$ takes values in sets; its $\kk$-linearization yields precisely
		the principal upset indicator $\kk[\uparrow q]$.  Dually, the corepresentable
		$\mathrm{Hom}_Q(-,q)$ linearizes to $\kk[\downarrow q]$.  For finite $Q$, these
		principal indicator modules behave like indecomposable projectives and injectives
		in the functor category, and they serve as building blocks for resolutions (compare
		with \cite{Mil20}).
	\end{remark}
	
	\subsection{1-parameter persistence}\label{sec:background:1d}
	
	This subsection recalls the classical one-parameter story, both to establish
	notation and to emphasize what breaks in the multiparameter setting.
	
	\subsubsection*{From filtrations to persistence modules}
	
	\begin{defn}[Filtration]\label{def:filtration}
		Let $Q$ be a poset.  A \emph{$Q$-filtration} of a topological space $X$ is a functor
		\[
		\mathcal{F}\colon Q\longrightarrow \mathbf{Top}
		\]
		such that $\mathcal{F}(q)$ is a subspace of $X$ for each $q$, and for $q\leq q'$ the
		morphism $\mathcal{F}(q\leq q')$ is the inclusion $\mathcal{F}(q)\hookrightarrow
		\mathcal{F}(q')$.  Equivalently, $\mathcal{F}$ is an order-preserving assignment of
		subspaces.
	\end{defn}
	
	Applying singular homology with coefficients in $\kk$ yields a persistence module.
	
	\begin{defn}[Persistent homology module]\label{def:persistent-homology-module}
		Given a $Q$-filtration $\mathcal{F}$ and an integer $i\geq 0$, the associated
		$i$th \emph{persistent homology module} is the composite functor
		\[
		H_i(\mathcal{F};\kk)\colon Q \xrightarrow{\ \mathcal{F}\ } \mathbf{Top}
		\xrightarrow{\ H_i(-;\kk)\ } \mathbf{Vect}_{\kk}.
		\]
	\end{defn}
	
	The case $Q=(\bbb{R},\leq)$ (or $Q=\bbb{Z}$) is the classical \emph{one-parameter}
	setting.  The filtration may arise, for example, from sublevel sets of a function
	$f\colon X\to \bbb{R}$, where $\mathcal{F}(t)=f^{-1}((-\infty,t])$.
	
	\subsubsection*{Interval modules and barcodes}
	
	The exceptional feature of one-parameter persistence is that (under mild finiteness
	hypotheses) modules admit a complete decomposition into interval summands.
	
	\begin{defn}[Interval module]\label{def:interval-module}
		Let $I\subseteq \bbb{R}$ be an interval (possibly infinite or half-open, depending on
		the chosen convention).  The \emph{interval module} $\kk_I$ is the $\bbb{R}$-module
		defined by
		\[
		\kk_I(t)=
		\begin{cases}
			\kk & t\in I,\\
			0 & t\notin I,
		\end{cases}
		\]
		with structure maps $\kk_I(s\leq t)$ equal to the identity on $\kk$ when $s,t\in I$
		and the zero map otherwise.
	\end{defn}
	
	To state the classification theorem, one usually imposes a finiteness condition,
	such as \emph{pointwise finite-dimensional} (p.f.d.), meaning $\dim_{\kk} M(t)<\infty$
	for all $t$, together with an appropriate tameness condition to rule out pathologies
	in the indexing by $\bbb{R}$ (see \cite{CB15,CDSGO16}).
	
	\begin{theorem}[Barcode decomposition, informal statement]\label{thm:barcode}
		Let $M$ be a p.f.d.\ persistence module indexed by $\bbb{R}$ satisfying a standard
		tameness condition (e.g.\ q-tameness).  Then $M$ decomposes as a direct sum of
		interval modules,
		\[
		M \cong \bigoplus_{\alpha\in A} \kk_{I_\alpha},
		\]
		and the multiset of intervals $\{I_\alpha\}$ is uniquely determined by $M$ up to
		reordering.  This multiset is the \emph{barcode} (or persistence diagram) of $M$.
	\end{theorem}
	
	The earliest formulations of this result for discrete indexing and finitely generated
	modules go back to persistent homology in computational topology; see, for example,
	\cite{ZC05,EH10}.  For $\bbb{R}$-indexed modules,
	a definitive representation-theoretic proof is given by Crawley-Boevey
	\cite{CB15}.
	
	\subsubsection*{Stability and distances}
	
	A major reason barcodes are useful in applications is \emph{stability}: small changes
	in input data yield small changes in the output barcode with respect to a natural
	metric.  There are several equivalent (or closely related) metrics in one parameter,
	notably the bottleneck distance on persistence diagrams and the interleaving distance
	on persistence modules; see \cite{CDSGO16,Les15} for
	systematic treatments.
	
	For the purposes of this thesis, we will use one-parameter stability primarily as a
	benchmark: it is an example of a situation where both (i) a complete discrete
	invariant exists (Theorem~\ref{thm:barcode}) and (ii) that invariant interacts well
	with perturbations.  Both features become substantially more complicated in multiple
	parameters.
	
	\subsection{Multiparameter persistence and why it is hard}\label{sec:background:multi-hard}
	
	\subsubsection*{Definition and basic examples}
	
	\begin{defn}[Multiparameter persistence module]\label{def:multi-parameter-module}
		Fix $n\geq 2$.  A (real-parameter) \emph{$n$-parameter persistence module} is a
		$\bbb{R}^n$-module, i.e.\ a functor $M\colon (\bbb{R}^n,\leq)\to\mathbf{Vect}_{\kk}$,
		where the order is coordinatewise:
		\[
		a=(a_1,\dots,a_n)\leq b=(b_1,\dots,b_n)\quad\Longleftrightarrow\quad
		a_i\leq b_i\ \text{for all }i.
		\]
		Likewise, a \emph{lattice} $n$-parameter persistence module is a $\bbb{Z}^n$-module.
	\end{defn}
	
	Such modules arise from \emph{multifiltrations} $\mathcal{F}\colon \bbb{R}^n\to
	\mathbf{Top}$, for instance from a function $f\colon X\to \bbb{R}^n$ via
	$\mathcal{F}(a)=f^{-1}((-\infty,a_1]\times\cdots\times(-\infty,a_n])$, and then
	applying homology as in Definition~\ref{def:persistent-homology-module}.
	
	\subsubsection*{Why barcodes do not generalize directly}
	
	When $n\geq 2$, the classification landscape changes qualitatively.  There is no
	analog of Theorem~\ref{thm:barcode} in which modules decompose uniquely into
	interval-like building blocks (rectangles, axis-aligned boxes, etc.).  One reason is
	representation-theoretic: the category of multiparameter persistence modules (even
	with finite-dimensionality restrictions) contains families of indecomposables
	depending on continuous parameters and exhibits behavior described as ``wild'' in
	the sense of quiver representation theory \cite{CZ09,Oud15}.
	
	At an informal level, ``wild'' means that the isomorphism classification problem is
	at least as complicated as classifying representations of arbitrary finite-dimensional
	algebras.  Consequently, one does not expect a complete discrete invariant such as a
	barcode to exist in general.
	
	\begin{remark}[What replaces barcodes?]\label{rmk:what-replaces-barcodes}
		In multiparameter persistence, one typically trades completeness for computability and
		stability: instead of a single complete invariant, one studies families of computable
		invariants (rank-based, homological, and/or slice-based) that are informative for
		specific tasks and data regimes.  Much of the design philosophy of TAMER-OP is to build
		a flexible pipeline that supports such families in a uniform algebraic framework.
	\end{remark}
	
	\subsubsection*{Invariants that still make sense}
	
	Even though complete classification fails, many useful invariants remain available.
	
	\begin{defn}[Rank invariant]\label{def:rank-invariant}
		Let $M$ be a $Q$-module.  The \emph{rank invariant} of $M$ is the function
		\[
		\rho_M \colon \{(q,q')\in Q\times Q : q\leq q'\}\longrightarrow \bbb{Z}_{\geq 0},
		\qquad
		\rho_M(q,q') := \mathrm{rank}\big(M(q\leq q')\big).
		\]
	\end{defn}
	
	\begin{defn}[Hilbert function / dimension function]\label{def:hilbert-function}
		The \emph{dimension function} (often called the \emph{Hilbert function} in graded
		settings) of a $Q$-module $M$ is
		\[
		h_M \colon Q\to \bbb{Z}_{\geq 0},\qquad h_M(q):=\dim_{\kk} M(q).
		\]
	\end{defn}
	
	\begin{remark}
		For $Q=\bbb{N}^n$ and suitably finite modules, $h_M$ is literally the Hilbert function
		of a multigraded module over a polynomial ring.  For general posets, it is still a
		basic and often stable summary of size.
	\end{remark}
	
	Another robust idea is to reduce to one parameter by restricting along monotone
	maps from a totally ordered set into $Q$.
	
	\begin{defn}[Restriction to a chain / slice]\label{def:slice}
		Let $C$ be a totally ordered set (often $C=\bbb{R}$ or $\bbb{Z}$) and let
		$\gamma\colon C\to Q$ be an order-preserving map.  The \emph{restriction} (or
		\emph{slice}) of a $Q$-module $M$ along $\gamma$ is the $C$-module $\gamma^*M$
		defined by composition:
		\[
		\gamma^*M := M\circ \gamma \colon C\to \mathbf{Vect}_{\kk}.
		\]
		When $C$ is totally ordered and $\gamma$ is injective, the image of $\gamma$ is a
		chain in $Q$.
	\end{defn}
	
	The restrictions $\gamma^*M$ are one-parameter persistence modules, hence admit
	barcodes under the usual hypotheses.  Slicing therefore packages multiparameter data
	into a family of one-parameter summaries; this idea is central in 2-parameter
	visualization tools such as RIVET \cite{RIVET}.
	
	Finally, there are homological invariants arising from resolutions and derived
	functors (Ext, Tor) that encode multi-parameter interactions.  These will be pursued
	systematically later in the thesis, but the guiding idea is already visible in
	Miller's program: introduce a finiteness condition strong enough to guarantee finite
	presentations and resolutions (and thus computable derived invariants), while weak
	enough to encompass common persistence modules arising from data \cite{Mil20}.
	
	\subsection{Two complementary viewpoints}\label{sec:background:viewpoints}
	
	Multiparameter persistence can be approached from at least two complementary
	perspectives.  The computational architecture of TAMER-OP is explicitly designed to
	move between them.
	
	\subsubsection*{Topological viewpoint: multifiltrations}
	
	At the source is topology: data are used to build a multifiltration
	$\mathcal{F}\colon \bbb{R}^n\to\mathbf{Top}$ (or $\bbb{Z}^n\to\mathbf{Top}$), and
	persistent homology applies a functor $H_i(-;\kk)$ to obtain
	\[
	M_i := H_i(\mathcal{F};\kk)\colon \bbb{R}^n\to \mathbf{Vect}_{\kk}.
	\]
	The topological viewpoint is advantageous for:
	\begin{itemize}
		\item motivating what invariants should mean geometrically (birth, death, and
		interactions across parameters);
		\item importing stability statements from the theory of interleavings and perturbation
		bounds; and
		\item making contact with geometric constructions of filtrations (sublevel sets, distance
		functions, Vietoris--Rips and \v{C}ech complexes, etc.; see \cite{EH10}).
	\end{itemize}
	
	\subsubsection*{Algebraic viewpoint: multigraded commutative algebra}
	
	When the indexing poset is discrete and additive, persistence modules are closely
	related to multigraded modules over polynomial rings.
	
	\begin{proposition}[Equivalence with multigraded modules (discrete setting)]
		\label{prop:graded-module-equivalence}
		Let $n\geq 1$ and consider the poset $(\bbb{N}^n,\leq)$ with coordinatewise order.
		Then the category of $\bbb{N}^n$-modules $M\colon \bbb{N}^n\to \mathbf{Vect}_{\kk}$
		is equivalent to the category of $\bbb{N}^n$-graded modules over the polynomial ring
		$\kk[x_1,\dots,x_n]$ with degree-preserving homomorphisms.
	\end{proposition}
	
	\begin{remark}[Sketch of correspondence]
		Given a persistence module $M$, form the graded $\kk$-vector space
		$\bigoplus_{a\in \bbb{N}^n} M(a)$ and let $x_i$ act by the structure maps
		$M(a\leq a+e_i)\colon M(a)\to M(a+e_i)$.  Functoriality ensures that the $x_i$ commute.
		Conversely, a graded $\kk[x_1,\dots,x_n]$-module determines the structure maps by
		multiplication with monomials.  See \cite{CZ09,Oud15} for
		detailed treatments in the persistence context.
	\end{remark}
	
	This algebraic viewpoint brings powerful tools---resolutions, Betti numbers, Hilbert
	series---but also reveals why classification becomes hard when
	$n\geq 2$: $\kk[x_1,\dots,x_n]$ is no longer a PID, and multigraded modules can have
	intricate syzygies and moduli.
	
	\subsubsection*{Finite posets and incidence algebras}
	
	TAMER-OP often works with finite posets (encoding posets) obtained by compressing a
	large or continuous parameter space.  For a finite poset $P$, $P$-modules can be
	viewed as finite-dimensional representations of the incidence algebra of $P$ (or of a
	quiver with commutativity relations determined by $P$).  This provides both:
	\begin{itemize}
		\item a representation-theoretic interpretation of projectives/injectives and
		resolutions (cf.\ Remark~\ref{rmk:yoneda-principal-upset}), and
		\item an algorithmic opportunity: for finite $P$, most constructions reduce to finite
		linear algebra, and many ``global'' constraints can be enforced by working only
		along cover relations (Hasse edges).
	\end{itemize}
	
	\subsection{Survey of partial solutions}
	\label{sec:background:partial-solutions}
	
	Because multiparameter persistence resists complete classification, a large body of
	work focuses on \emph{partial} invariants and practical pipelines.  The following
	four themes recur throughout the literature and motivate the design choices of this
	thesis.
	
	\subsubsection*{(i) Slicing and fibered barcodes}
	
	The restriction construction in Definition~\ref{def:slice} reduces a multiparameter
	module to a one-parameter module, hence to a barcode.  Varying the slice (for
	instance, varying a line of positive slope in $\bbb{R}^2$) yields a family of barcodes
	that can be used for visualization or for defining distances by integrating (or
	taking suprema of) one-parameter distances.  In the 2-parameter case, the
	\emph{fibered barcode} viewpoint underlies the RIVET framework \cite{RIVET}.
	
	From a computational standpoint, slicing has two key advantages:
	\begin{enumerate}[label=(\roman*)]
		\item it leverages mature one-parameter algorithms, and
		\item it provides stable summaries for applications even when the full module is too
		large to manipulate directly.
	\end{enumerate}
	Its main limitation is incompleteness: different modules can have identical slices.
	
	\subsubsection*{(ii) Rank-based invariants and M\"obius inversion}
	
	The rank invariant (Definition~\ref{def:rank-invariant}) is one of the most basic
	multiparameter invariants.  In one parameter, the rank function determines the
	barcode, and the barcode can be recovered from ranks by an inclusion--exclusion
	(M\"obius inversion) formula.  In higher parameters, an analogous inversion yields
	\emph{signed} measures on axis-aligned rectangles (sometimes called \emph{signed
		barcodes} or \emph{rectangle measures}).  These objects can be used as:
	\begin{itemize}
		\item sparse, computable summaries of the rank invariant on a chosen grid, and
		\item feature representations for statistics/learning pipelines (e.g.\ kernels on
		signed measures).
	\end{itemize}
	From the algebraic side, such M\"obius inversions are naturally phrased on products of
	chains, and they fit well with finite encodings (cf.\ (iv) below).
	
	\subsubsection*{(iii) Commutative algebra invariants via resolutions}
	
	Under the equivalence in Proposition~\ref{prop:graded-module-equivalence}, a
	$\bbb{N}^n$-module becomes a multigraded $\kk[x_1,\dots,x_n]$-module, and one can
	study:
	\begin{itemize}
		\item minimal free resolutions and multigraded Betti numbers,
		\item Tor/Ext groups and their algebra structures,
		\item associated primes and socles, and
		\item Hilbert series and local cohomology.
	\end{itemize}
	These invariants detect multi-parameter interactions that rank functions and slices
	can miss.  However, they require a finiteness condition strong enough to ensure that
	resolutions are finite in a computational sense.  For real-parameter modules, the
	usual noetherian hypotheses fail, motivating Miller's theory of \emph{tameness} as a
	replacement \cite{Mil20}.
	
	\subsubsection*{(iv) Geometric/combinatorial encoding strategies}
	
	A guiding principle in this thesis is that many modules of interest are
	\emph{piecewise constant} over the parameter space in a controlled way.  An
	\emph{encoding} compresses the indexing poset $Q$ (often $\bbb{R}^n$ or $\bbb{Z}^n$)
	to a finite poset $P$ so that the module factors through $P$ and all subsequent
	computations may be carried out on the finite object.
	
	In Miller's framework, tameness can be characterized in several equivalent ways,
	including (among others) the existence of a finite encoding and the existence of
	certain finite ``fringe'' presentations by upset and downset indicator modules
	\cite{Mil20}.  The computational moral is that once such a finite encoding
	is constructed, one can safely deploy the abelian and homological algebra of finite
	poset modules (kernels, images, Ext/Tor, resolutions) using finite linear algebra.
	
	\subsection{Brief software landscape}
	\label{sec:background:software-landscape}
	
	This thesis is not primarily a software survey, but it is useful to clarify how
	different toolchains emphasize different parts of the multiparameter landscape.
	We will return to more detailed comparisons in later sections, after the TAMER-OP data
	model and pipeline are established.
	
	\begin{itemize}
		\item \textbf{RIVET.} In the 2-parameter case, RIVET is centered on the fibered
		barcode: it supports interactive exploration of how one-parameter barcodes vary over
		a family of lines, and it provides infrastructure for rank-based visualization and
		related distances \cite{RIVET}.
		
		\item \textbf{multipers.} The multipers ecosystem (as commonly used in applied TDA)
		targets scalable computation of multiparameter summaries and feature maps (often
		slice-based, rank-based, or image/landscape-type constructions).  Its emphasis is
		typically on practical pipelines rather than complete algebraic structure.  (A precise
		feature-by-feature comparison depends on the chosen backends and will be addressed
		later.)
		
		\item \textbf{TAMER-OP.} TAMER-OP is designed around Miller's tame-module viewpoint:
		build or ingest a presentation over a large indexing poset, compress it via a finite
		encoding poset, and then compute a broad family of invariants (rank-based,
		homological/derived, geometric, and statistical) by systematically exploiting abelian
		and homological algebra over the finite encoding.
	\end{itemize}
	
	\begin{remark}[Setting expectations for comparisons]\label{rmk:comparison-expectations}
		Comparisons are most meaningful when the underlying mathematical object is held fixed.
		For example, comparing fibered barcodes is natural when both toolchains work with
		2-parameter modules and slice-based summaries.  Comparing resolution-based invariants
		is meaningful when both toolchains expose (or can emulate) multigraded homological
		algebra.  TAMER-OP's main distinguishing feature is that it treats \emph{finite encodings}
		as first-class objects, so that many invariants can be computed uniformly once the
		encoding is in place.
	\end{remark}
	
	\subsection{Notation and conventions}\label{sec:background:conventions}
	
	\subsubsection*{Coefficients and exactness}
	
	Unless stated otherwise, modules take values in $\mathbf{Vect}_{\kk}$ for a fixed
	field $\kk$.  In implementations, there is a persistent tradeoff between speed and
	certifiability:
	\begin{itemize}
		\item Over floating-point fields (e.g.\ $\bbb{R}$), numerical linear algebra is fast
		but can introduce rank errors near degeneracies.
		\item Over exact fields (notably $\bbb{Q}$), ranks, kernels, and Smith/row-reduction
		type computations are certifiable but may be more expensive.
	\end{itemize}
	TAMER-OP is designed so that exact computations over $\bbb{Q}$ can be used when needed,
	especially for invariants derived from ranks and homological algebra (kernels,
	cokernels, Ext/Tor dimensions).
	
	\subsubsection*{Orders on $\bbb{Z}^n$ and $\bbb{R}^n$}
	
	We use the coordinatewise order on $\bbb{Z}^n$ and $\bbb{R}^n$:
	\[
	a\leq b \iff a_i\leq b_i\ \text{for all }i.
	\]
	When writing $a<b$, we mean $a\leq b$ and $a\neq b$.  The standard basis vectors are
	$e_1,\dots,e_n$.
	
	For a finite poset $P$, we generally denote its order by $\leq$ as well; the intended
	meaning should be clear from context.  When $P$ is used as an encoding poset for an
	ambient parameter space, elements of $P$ are sometimes called \emph{regions}.
	
	\subsubsection*{Birth/death terminology and (co)presentations}
	
	When working with indicator modules (Definition~\ref{def:indicator-modules}) and their
	linear combinations, we follow the ``birth/death'' intuition from persistence:
	\begin{itemize}
		\item Upsets model \emph{birth} regions (where generators become present and persist
		upward in the order).
		\item Downsets model \emph{death} constraints (where relations become enforced downward
		in the order).
	\end{itemize}
	A \emph{presentation} will typically mean a cokernel of a map between free objects
	built from upset indicators, while a \emph{copresentation} dually involves downset
	indicators.  In Miller's terminology, \emph{fringe presentations} mix these two
	directions and are tailored to tame modules over general posets \cite{Mil20}.
	
	% ============================================================
	% Section 3: Tame modules
	% ============================================================
	
	\section{Tame modules}
	\label{sec:tame-modules}
	
	Multiparameter persistence is often presented as a story about functoriality: a
	multifiltration of spaces (or chain complexes) gives rise to a functor from a parameter
	poset into vector spaces, and this functor is the algebraic object one studies.
	The difficulty is that, once the parameter poset has dimension $\ge 2$, there is no
	reasonable discrete classification of all such functors, even under strong finiteness
	conditions such as pointwise finite-dimensionality; in representation-theoretic language,
	the classification problem becomes ``wild'' in general \cite{CZ09}.
	For computational purposes one therefore needs a \emph{hypothesis of controlled
		complexity}---a notion of finiteness that is strong enough to guarantee that modules admit
	finite descriptions and are stable under the constructions one cares about, but flexible
	enough to encompass the examples arising in applications.
	
	In this thesis we adopt Ezra Miller's framework of \emph{tame modules over posets}
	\cite{Mil20}.  The guiding idea is that the module should have only finitely many
	distinct ``regions of algebraic behavior'' in parameter space, so that it can be
	represented by \emph{finite data} after passing through a suitable \emph{encoding}.
	A primary goal of TAMER-OP is to turn this philosophy into concrete data structures and
	algorithms; the present section sets up the mathematical backbone that makes such an
	implementation sensible.
	
	\subsection{Why ``tame'' is the right hypothesis for computation}
	\label{subsec:why-tame}
	
	The classical $1$-parameter theory enjoys a powerful structure theorem:
	under mild hypotheses, persistence modules decompose as direct sums of interval modules,
	and the resulting multiset of intervals is encoded by a barcode or persistence diagram.
	This viewpoint underlies much of applied persistent homology, as well as practical
	algorithms for computing invariants from filtered data \cite{ZC05}.
	One rigorous formulation of this classification is due to Crawley--Boevey:
	pointwise finite-dimensional persistence modules over a totally ordered index set admit
	interval decompositions under appropriate finiteness hypotheses \cite{CB15}.
	
	The multiparameter setting behaves fundamentally differently.  Even for
	$\bbb{Z}^n$-indexed modules with $n\ge 2$, indecomposables occur in families and the
	category does not admit a discrete complete invariant analogous to a barcode \cite{CZ09}.
	From a computational perspective, this means that one should not expect a \emph{complete}
	classification of multiparameter persistence modules by a finite combinatorial object.
	Instead, one aims for a framework in which
	\begin{enumerate}[label=(\roman*)]
		\item the objects of interest admit \emph{finite representations},
		\item the relevant algebraic constructions (kernels/cokernels, derived functors, etc.)
		stay within the class, and
		\item the invariants one computes are \emph{functorial} and stable under natural
		operations on the input data.
	\end{enumerate}
	
	Miller's notion of tameness is designed with exactly these desiderata in mind.
	Very roughly, a module is tame if it is constant on finitely many regions of parameter
	space and therefore can be expressed via a finite ``model poset'' that records the
	comparability relations among these regions \cite{Mil20}.
	The definition is arranged so that the category of tame modules (with an appropriate
	restriction on morphisms) is abelian, hence supports homological algebra; moreover,
	tameness admits several equivalent characterizations---combinatorial, topological, and
	homological---that are captured by the syzygy theorem (Theorem~\ref{thm:syzygy} below).
	
	\subsection{Encoding posets and encoding maps}
	\label{subsec:encoding-posets}
	
	The central finiteness notion used throughout this thesis is \emph{finite poset encoding}.
	Intuitively, an encoding replaces a large (typically infinite) poset $Q$ by a smaller
	poset $P$ whose elements represent ``regions'' of the parameter space, such that the
	module on $Q$ is obtained by pullback from a module on $P$.
	
	\begin{defn}[Poset encoding {\cite[Def.~4.1]{Mil20}}]
		\label{def:poset-encoding}
		Let $Q$ be a poset, viewed as a category in the usual way.
		An \emph{encoding} of a $Q$-module $M$ by a poset $P$ consists of
		\begin{enumerate}[label=(\roman*)]
			\item an order-preserving map (equivalently, a functor) $\pi \colon Q \to P$, and
			\item a $P$-module $H \colon P \to \mathbf{Vect}_{\kk}$
		\end{enumerate}
		together with an isomorphism of $Q$-modules
		\[
		M \ \cong \ \pi^*H \ := \ H\circ \pi .
		\]
		The encoding is called \emph{finite} if $P$ has finitely many elements and each stalk
		$H(p)$ is a finite-dimensional $\kk$-vector space.
	\end{defn}
	
	\begin{remark}[Encodings as ``finite models'']
		\label{rem:encoding-fibers}
		If $M\cong \pi^*H$, then $M$ is constant on each fiber $\pi^{-1}(p)$: for any
		$q,q'\in Q$ with $\pi(q)=\pi(q')=p$, the only morphism $p\to p$ in the poset category
		is the identity, so the structure maps of $M$ restrict to identities on the shared stalk
		$H(p)$.  Thus a finite encoding exhibits $M$ as having finitely many ``constant
		regions'' indexed by $P$.
		Conversely, in many geometric situations a partition of the parameter space into finitely
		many constant regions determines a finite encoding by taking $P$ to be the poset of
		regions ordered by reachability (or by closure relations); see \cite{Mil20} for the
		general equivalence among the various constructions.
	\end{remark}
	
	The functorial viewpoint is especially important in what follows, so we record the basic
	formalism.
	
	\begin{defn}[Pullback along a poset morphism]
		\label{def:pullback}
		Let $\pi\colon Q\to P$ be a poset morphism.  Precomposition defines a functor
		\[
		\pi^* \colon \mathbf{Vect}_{\kk}^{P} \longrightarrow \mathbf{Vect}_{\kk}^{Q},
		\qquad
		H \longmapsto H\circ \pi.
		\]
		We refer to $\pi^*$ as the \emph{pullback} (or restriction) functor along $\pi$.
	\end{defn}
	
	\begin{lemma}[Exactness of pullback]
		\label{lem:pullback-exact}
		For any poset morphism $\pi\colon Q\to P$, the pullback functor $\pi^*$ is exact.
		In particular, it preserves finite limits and colimits (kernels, cokernels, and finite
		direct sums).
	\end{lemma}
	
	\begin{proof}
		The category $\mathbf{Vect}_{\kk}^{P}$ of $P$-modules is a functor category with values
		in $\mathbf{Vect}_{\kk}$, hence limits and colimits are computed pointwise.
		Given a short exact sequence $0\to H'\to H\to H''\to 0$ of $P$-modules, evaluating at
		any $p\in P$ yields a short exact sequence of vector spaces.  Precomposing with $\pi$
		and then evaluating at $q\in Q$ simply evaluates the original sequence at $\pi(q)$, so
		exactness is preserved at every stalk, hence globally.
	\end{proof}
	
	Definition~\ref{def:poset-encoding} formalizes the idea that the module ``factors
	through'' a finite poset.  The following terminology, also standard in \cite{Mil20}, is
	useful when one wishes to emphasize the dependence on a fixed encoding map.
	
	\begin{defn}[Subordinate encoding {\cite[Def.~4.7]{Mil20}}]
		\label{def:subordinate}
		An encoding map $\pi\colon Q\to P$ is said to be \emph{subordinate} to a $Q$-module $M$
		if there exists a $P$-module $H$ such that $M\cong \pi^*H$.
	\end{defn}
	
	\begin{example}[Indicator modules admit trivial finite encodings {\cite[Ex.~4.6]{Mil20}}]
		\label{ex:indicator-encoding}
		Let $U\subseteq Q$ be an upset (upward-closed set).  Then the indicator module $\kk\{U\}$
		(defined in Subsection~\ref{subsec:fringe-presentations}) admits a finite encoding by the
		two-element chain $0<1$ via the map $\pi\colon Q\to\{0<1\}$ sending $U$ to $1$ and
		$Q\setminus U$ to $0$; the corresponding module $H$ has $H(1)=\kk$ and $H(0)=0$.
		Dually, any downset indicator module $\kk[D]$ is encoded by the same chain with the
		roles of $0$ and $1$ reversed.
	\end{example}
	
	\subsection{Definitions of tameness in this framework}
	\label{subsec:def-tameness}
	
	The previous subsection defined what it means for a module to be encoded by a fixed
	poset $P$.  Tameness is the assertion that such an encoding exists with $P$ finite.
	
	\begin{defn}[$P$-tame and tame modules]
		\label{def:tame-module}
		Fix a poset $Q$.
		\begin{enumerate}[label=(\roman*)]
			\item Let $P$ be a finite poset.  A $Q$-module $M$ is \emph{$P$-tame} if there exists a
			poset morphism $\pi\colon Q\to P$ and a $P$-module $H$ with finite-dimensional stalks
			such that $M\cong \pi^*H$.
			\item A $Q$-module $M$ is \emph{tame} if it is $P$-tame for some finite poset $P$.
			Equivalently, $M$ admits a finite poset encoding subordinate to it in the sense of
			Definition~\ref{def:poset-encoding}.
		\end{enumerate}
	\end{defn}
	
	\begin{remark}[Comparison with other finiteness notions]
		\label{rem:other-tameness}
		There are several finiteness conditions in the persistence literature, especially over
		$\bbb{R}^n$ and $\bbb{Z}^n$, including variants of ``constructibility'' and ``$q$-tameness''
		based on finiteness of ranks of structure maps.
		Those conditions are well adapted to stability questions and to certain metric structures,
		but they do not by themselves supply a finite combinatorial model that supports the full
		suite of homological computations one might wish to perform.
		Finite encodings are a \emph{stronger, more combinatorial} hypothesis: they guarantee
		that all information about $M$ is carried by a finite poset $P$ together with finite
		linear algebra on the stalks of a $P$-module.  In particular, tameness is expressly
		designed to enable algorithmic computation in principle \cite{Mil20}.
	\end{remark}
	
	A subtle point is that \emph{not every} $Q$-module homomorphism between tame modules
	should be regarded as ``computably tame''.  Allowing all homomorphisms can take one
	outside the tame world (for example, kernels of arbitrary maps need not be finitely
	encoded).  The correct remedy is to restrict to morphisms compatible with a common
	finite description.
	
	\begin{defn}[Tame morphisms and the tame category {\cite[Def.~4.27, Def.~4.29]{Mil20}}]
		\label{def:tame-morphism}
		Let $M,N$ be tame $Q$-modules.
		A homomorphism $\varphi\colon M\to N$ of $Q$-modules is called \emph{tame} if there
		exists a finite ``constant subdivision'' of $Q$ that is subordinate to both $M$ and $N$,
		such that on each constant region the induced linear map is independent of the chosen
		parameter value (see \cite[Def.~4.27]{Mil20} for the precise formulation).
		The \emph{category of tame $Q$-modules} is the category whose objects are tame
		$Q$-modules and whose morphisms are tame homomorphisms.
	\end{defn}
	
	\begin{theorem}[Tame modules form an abelian category {\cite[Prop.~4.31]{Mil20}}]
		\label{thm:tame-abelian}
		For any poset $Q$, the category of tame $Q$-modules (with tame morphisms) is abelian.
		In particular, it has kernels and cokernels, finite products and coproducts, and every
		monomorphism (resp.\ epimorphism) is a kernel (resp.\ cokernel).
	\end{theorem}
	
	\subsection{Why non-abelianity fears are usually irrelevant in computation}
	\label{subsec:tame-nonabelianity-practice}

	The worry voiced by many users of multiparameter persistence is the following.
	The ambient category $\mathbf{Vect}_{\kk}^{Q}$ of $Q$-modules is abelian, so it is natural
	to form kernels, cokernels, homology, and derived functors there.
	However, the full subcollection of \emph{tame} objects inside $\mathbf{Vect}_{\kk}^{Q}$ is
	not automatically stable under these operations if one allows \emph{arbitrary} natural
	transformations between tame objects.  In other words, tameness can fail to be preserved
	if one performs abelian constructions in $\mathbf{Vect}_{\kk}^{Q}$ without restricting the
	class of morphisms.

	This subsection unpacks the mathematical content behind that concern and explains why,
	from the standpoint of computation, the bad behavior is almost never encountered.
	The key point is that there are two distinct notions of ``morphism between tame modules'',
	and only one of them is compatible with finite data structures.

	\subsubsection*{Homomorphisms of tame objects versus tame morphisms}

	Recall from Remark~\ref{rem:tame-morphisms-practice} that a finite encoding
	$\pi\colon Q\to P$ converts a tame $Q$-module into a $P$-module for a finite poset $P$.
	In that finite setting, morphisms are ordinary natural transformations in
	$\mathbf{Vect}_{\kk}^{P}$, hence are controlled by finitely many linear maps.

	Inside $\mathbf{Vect}_{\kk}^{Q}$, however, a natural transformation
	$\varphi\colon M\to N$ between tame modules can vary on infinitely many degrees even when
	$M$ and $N$ are each constant on finitely many regions.  Miller isolates the maps that
	are compatible with finite encodings by requiring \emph{constancy of the transformation
	on constant regions}:
	a homomorphism $\varphi$ is called \emph{tame} if there exists a finite constant subdivision
	subordinate to both $M$ and $N$ such that on each region $I$ the composite map
	$M_I\to M_q\to N_q\to N_I$ is independent of $q\in I$ \cite[Def.~4.27]{Mil20}.
	The category of tame modules uses these tame homomorphisms as its morphisms
	\cite[Def.~4.29 and Rem.~4.30]{Mil20}.  The resulting category is abelian by
	Theorem~\ref{thm:tame-abelian}.

	\subsubsection*{A concrete failure mode}

	The distinction between arbitrary homomorphisms and tame morphisms is not cosmetic.
	The following example demonstrates that kernels can leave the tame world if the morphism
	is allowed to vary infinitely within a constant region.

	\begin{example}[A kernel of a homomorphism of tame modules need not be tame {\cite[Ex.~4.25]{Mil20}}]
		\label{ex:kernel-not-tame}
		Let $Q=\mathbb{R}^2$ with the coordinatewise partial order, and let
		$L=\{(x,y)\in\mathbb{R}^2 : x+y=1\}$ be the antidiagonal line.
		Let $U\subseteq Q$ be the closed upset above $L$, and let $U^{\circ}$ be its interior.
		Consider the tame modules
		\[
			M=\kk\{U\}\oplus \kk\{U\},
			\qquad
			N=\kk\{U\}/\kk[U^{\circ}],
		\]
		where $N$ is supported on $L$.
		Define a surjection $\varphi\colon M\twoheadrightarrow N$ by taking $\varphi$ to be
		zero on $U^{\circ}$ and, for each point $(x,y)\in L$, sending the two standard basis
		vectors of $M_{(x,y)}=\kk^2$ to the scalars $y$ and $-x$ in $N_{(x,y)}=\kk$.

		Then $K=\ker\varphi$ is \emph{not} tame.  Indeed, for $(x,y)\in L$, the stalk $K_{(x,y)}$
		is the one-dimensional subspace of $\kk^2$ spanned by $(x,y)$, and these lines are pairwise
		distinct as $(x,y)$ varies along $L$.  Consequently, any constant subdivision subordinate to $K$
		must separate distinct points of $L$ into distinct constant regions, forcing infinitely many regions.
	\end{example}

	The phenomenon in Example~\ref{ex:kernel-not-tame} is that $M$ and $N$ are each constant on a
	finite subdivision of $Q$ (three regions: below $L$, on $L$, and above $L$), but the map
	$\varphi$ is \emph{not} constant on the region $L$ itself.  Hence $\varphi$ is not a tame
	homomorphism, so there is no reason for its kernel to remain tame.

	\subsubsection*{The tame category is closed under the usual constructions}

	Once morphisms are required to be tame, the basic abelian operations are safe.

	\begin{lemma}[Kernels and cokernels of tame homomorphisms are tame {\cite[Lem.~4.28]{Mil20}}]
		\label{lem:tame-kernel-cokernel}
		Let $\varphi\colon M\to N$ be a tame homomorphism of $Q$-modules.
		Then $\ker\varphi$ and $\coker\varphi$ are tame modules, and the canonical maps
		$\ker\varphi\to M$ and $N\to\coker\varphi$ are tame morphisms.
		The same statement holds in the semialgebraic, PL, subanalytic, or class~$X$ settings
		when the constant subdivision witnessing tameness is of the corresponding type.
	\end{lemma}

	\begin{proof}
		Choose a finite constant subdivision of $Q$ that is subordinate to $M$, $N$, and $\varphi$
		(as required in the definition of a tame homomorphism).
		For each constant region $I$, fix identifications $M_q\cong M_I$ and $N_q\cong N_I$
		for all $q\in I$ so that the structure maps of $M$ and $N$ agree with the transition
		maps among the regionwise vector spaces.
		By tameness of $\varphi$, the induced linear map $M_I\to N_I$ is independent of $q\in I$.

		Define regionwise vector spaces
		\[
			K_I:=\ker(M_I\to N_I),
			\qquad
			C_I:=\coker(M_I\to N_I).
		\]
		Using the structure maps between regions, these assemble into $Q$-modules $K$ and $C$.
		By construction, $K=\ker\varphi$ and $C=\coker\varphi$ in the abelian category
		$\mathbf{Vect}_{\kk}^{Q}$, and the chosen constant subdivision is subordinate to both $K$
		and $C$.  Hence $K$ and $C$ are tame.  The auxiliary hypotheses (PL, semialgebraic, etc.)
		are preserved because the same subdivision witnesses tameness in those settings.
	\end{proof}

	Lemma~\ref{lem:tame-kernel-cokernel} is the formal version of the slogan that
	``computations with tame modules stay tame'' as long as the maps used in the computation
	come from finite data.  Example~\ref{ex:kernel-not-tame} shows that if one insists on allowing
	\emph{all} homomorphisms between tame objects, then kernels and cokernels can escape the tame
	world, but those escaping morphisms are precisely the ones that do not admit a finite encoding.

	\subsubsection*{Why the pathological maps do not appear in algorithms}

	Any explicit computation in this thesis ultimately proceeds by passing to some finite encoding
	$\pi\colon Q\to P$ and working in the finite functor category $\mathbf{Vect}_{\kk}^{P}$.
	In that regime, every morphism is represented by finitely many linear maps, and pulling back
	along $\pi$ produces a tame homomorphism in the sense of Definition~\ref{def:tame-morphism}.
	Exactness of pullback (Lemma~\ref{lem:pullback-exact}) then implies that kernels and cokernels
	computed at the finite level remain finite after pulling back.

	More explicitly, if $M\cong\pi^{*}H$ and $N\cong\pi^{*}G$ for $P$-modules $H$ and $G$, and if
	$\bar\varphi\colon H\to G$ is a morphism in $\mathbf{Vect}_{\kk}^{P}$, then
	$\varphi:=\pi^{*}\bar\varphi$ is a tame homomorphism of $Q$-modules, and
	\[
		\ker(\varphi)\cong \pi^{*}(\ker\bar\varphi),
		\qquad
		\coker(\varphi)\cong \pi^{*}(\coker\bar\varphi),
	\]
	because $\pi^{*}$ is exact.  Therefore the outputs of the standard constructions
	(kernel, cokernel, homology of a complex, and derived functors computed via finite
	resolutions) remain tame whenever the computation is expressed at the level of a finite
	encoding.  In particular, the ``risky'' behavior in Example~\ref{ex:kernel-not-tame} cannot
	occur unless one is willing to specify a morphism that varies on infinitely many degrees
	inside a constant region, which is not compatible with the finite data structures used for
	encoding and computation.

	\subsubsection*{Recent results that restore full abelianity under geometric hypotheses}

	The preceding discussion explains why Miller restricts morphisms in the tame category.
	There is also a complementary line of results showing that, for many geometric classes of
	finite encodings, the restriction is in practice unnecessary: under suitable hypotheses on
	encoding fibers, every homomorphism between encodable modules can be realized inside a
	finite encoding.

	One precise formulation is due to Waas \cite{Waa24}, who introduces a connectivity condition
	on encoding fibers ensuring that any finite set of encodable modules admits a \emph{common}
	encoding that captures \emph{all} morphisms between them \cite[Prop.~3.5]{Waa24}.
	In particular, the resulting class of encodable persistence modules forms a full abelian
	subcategory of $\mathbf{Vect}_{\kk}^{Q}$ \cite[Thm.~1.1]{Waa24}.
	For parameter posets $Q=\mathbb{R}^n$, his hypotheses are satisfied by natural families of
	``geometric'' regions built from topologically closed PL, semialgebraic, or subanalytic
	upsets (and their complements and finite unions).

	Conceptually, the point is that the failure in Example~\ref{ex:kernel-not-tame} relies on
	large disconnected families of degrees (here the continuum of points on $L$) along which a map
	can vary independently.  Connectivity hypotheses on encoding fibers rule out that degree of
	freedom, forcing morphisms to be controlled by finite combinatorial data.  Waas also notes that
	one can alternatively prove the corresponding subanalytic statement by translating the problem
	into constructible sheaf theory and invoking results of Miller and Kashiwara-Schapira
	\cite[Rem.~3.19]{Waa24}, \cite{Mil23,KS18}.

	These developments give a precise mathematical interpretation of the informal claim that
	``if the input is tame, then the output is almost always tame'': the non-tame outputs arise
	from morphisms that do not admit finite encodings, and such morphisms are excluded either
	(i) by working in the tame category (tame morphisms by definition), or
	(ii) by restricting attention to the geometric subclasses of tame modules for which all
	relevant morphisms are automatically encodable.

	\subsection{Fringe presentations as a computable interface}
	\label{subsec:fringe-presentations}
	
	Finite encodings are the most conceptually direct way to impose combinatorial finiteness,
	but they are not always the most convenient format for \emph{input} or for certain
	computations.  Miller's framework supplies a complementary algebraic interface:
	\emph{fringe presentations}.  These generalize the idea of presenting a module via
	generators and relations, where generators are organized as upsets (birth regions) and
	relations as downsets (death regions).
	
	We begin by recalling the relevant combinatorics of subsets of a poset.
	
	\begin{defn}[Upsets and downsets]
		\label{def:upset-downset}
		Let $Q$ be a poset.
		\begin{enumerate}[label=(\roman*)]
			\item An \emph{upset} (or upper set) in $Q$ is a subset $U\subseteq Q$ such that
			$q\in U$ and $q\le q'$ imply $q'\in U$.
			\item A \emph{downset} (or lower set) in $Q$ is a subset $D\subseteq Q$ such that
			$q\in D$ and $q'\le q$ imply $q'\in D$.
		\end{enumerate}
		For $q\in Q$, the \emph{principal upset} generated by $q$ is $\uparrow q := \{q'\in Q : q\le q'\}$,
		and the \emph{principal downset} is $\downarrow q := \{q'\in Q : q'\le q\}$.
	\end{defn}
	
	\begin{defn}[Indicator modules for upsets and downsets]
		\label{def:indicator-modules}
		Let $Q$ be a poset and let $S\subseteq Q$.
		The \emph{indicator module} (or \emph{indicator $Q$-module}) $\kk[S]$ is the
		$Q$-module defined on objects by
		\[
		\kk[S](q) :=
		\begin{cases}
			\kk & \text{if } q\in S,\\
			0   & \text{if } q\notin S,
		\end{cases}
		\]
		and on morphisms by the rule that for $q\le q'$ the structure map
		$\kk[S](q\le q')\colon \kk[S](q)\to \kk[S](q')$ is the identity $\kk\to\kk$ when both
		$q,q'\in S$, and is $0$ otherwise.
		When $S=U$ is an upset, $\kk[U]$ is called an \emph{upset module}; when $S=D$ is a
		downset, $\kk[D]$ is called a \emph{downset module}.
	\end{defn}
	
	\begin{remark}[Projectives and injectives on a finite poset]
		\label{rem:proj-inj-finite-poset}
		When $P$ is a finite poset, the category $\mathbf{Vect}_{\kk}^{P}$ can be identified
		with the category of finite-dimensional representations of the Hasse quiver of $P$ with
		commutativity relations (equivalently, with modules over the incidence algebra of $P$).
		In this setting, principal upset modules $\kk[\uparrow p]$ are the indecomposable
		projective objects and principal downset modules $\kk[\downarrow p]$ are the
		indecomposable injective objects.  This is one reason indicator modules are the correct
		building blocks for computable homological algebra; see \cite{Mil20} for a detailed
		development in the general poset setting.
	\end{remark}
	
	The next definition packages a tame module using finitely many upsets and downsets
	together with a matrix over $\kk$.  This is the abstraction of ``birth'' and ``death''
	data familiar from barcodes, but adapted to general posets.
	
	\begin{defn}[Fringe presentation {\cite[Def.~3.16]{Mil20}}]
		\label{def:fringe-presentation}
		Let $Q$ be a poset and let $M$ be a $Q$-module.
		A \emph{fringe presentation} of $M$ consists of
		\begin{enumerate}[label=(\roman*)]
			\item a direct sum $F = \bigoplus_{i\in I}\kk[U_i]$ of upset modules,
			\item a direct sum $E = \bigoplus_{j\in J}\kk[D_j]$ of downset modules, and
			\item a $Q$-module homomorphism $\varphi\colon F\to E$
		\end{enumerate}
		such that $\im(\varphi)\cong M$ and the component maps
		$\kk[U_i]\to \kk[D_j]$ satisfy a mild connectivity condition (see \cite[Def.~3.14]{Mil20}).
		The fringe presentation is \emph{finite} if the index sets $I$ and $J$ are finite.
	\end{defn}
	
	\begin{remark}[Relation to classical presentations]
		\label{rem:fringe-as-presentation}
		If $Q=\bbb{Z}^n$ with the usual partial order, then $Q$-modules can be interpreted as
		$\bbb{Z}^n$-graded modules over the polynomial ring $\kk[x_1,\dots,x_n]$ (Section~2.4).
		In that language, upset modules behave like ``free/flat'' pieces and downset modules
		behave like ``injective'' pieces; the map $\varphi$ organizes generators (birth upsets)
		and relations (death downsets) in a way that is robust under changes of geometry and
		does not depend on choosing a coordinate-aligned region decomposition.
	\end{remark}
	
	Fringe presentations become especially concrete once one writes the map $\varphi$ in
	coordinates.
	
	\begin{defn}[Monomial matrix {\cite[Def.~3.17]{Mil20}}]
		\label{def:monomial-matrix}
		Let $\varphi\colon F\to E$ be a \emph{finite} fringe presentation with
		$F=\bigoplus_{i=1}^{m}\kk[U_i]$ and $E=\bigoplus_{j=1}^{r}\kk[D_j]$.
		Choosing the natural bases on the $1$-dimensional stalks of each indicator summand,
		the map $\varphi$ can be encoded by an $r\times m$ matrix $\Phi=(\Phi_{j i})$ with
		entries in $\kk$, whose columns are indexed by the birth upsets $U_i$ and whose rows
		are indexed by the death downsets $D_j$.  We refer to $\Phi$ as a \emph{monomial matrix}
		for $\varphi$.
	\end{defn}
	
	\begin{lemma}[Support condition for monomial matrices {\cite[Prop.~3.18]{Mil20}}]
		\label{lem:monomial-support}
		With notation as in Definition~\ref{def:monomial-matrix}, one necessarily has
		$\Phi_{j i}=0$ whenever $U_i\cap D_j=\varnothing$.  Conversely, any matrix
		$\Phi\in \kk^{r\times m}$ satisfying $\Phi_{j i}=0$ when $U_i\cap D_j=\varnothing$
		determines a fringe presentation map $\varphi\colon F\to E$.
	\end{lemma}
	
	\begin{example}[A single interval in one parameter]
		\label{ex:interval-fringe}
		Let $Q=\bbb{R}$ with its total order and consider the interval module $\kk[a,b)$.
		This module can be realized as the image of a map from the upset module $\kk[a,\infty)$
		to the downset module $\kk(-\infty,b)$, and the corresponding monomial matrix is the
		$1\times 1$ matrix $[1]$.  For a direct sum of intervals (a barcode), the matching of
		births to deaths is reflected by a monomial matrix that is a permutation matrix after
		choosing orderings of births and deaths; compare \cite[Ex.~3.19]{Mil20} and the
		interval decomposition theorem \cite{CB15}.
	\end{example}
	
	\subsection{What tameness buys you: mathematical consequences}
	\label{subsec:tame-math}
	
	The principal reason tameness is a useful organizing principle is that it is not a single
	isolated definition but rather part of a web of equivalent finiteness conditions, each
	suited to different tasks (geometric modeling, algebraic presentation, homological
	computation).  The main structural result is Miller's syzygy theorem.
	
	\begin{theorem}[Syzygy theorem {\cite[Thm.~6.12]{Mil20}}]
		\label{thm:syzygy}
		Let $Q$ be a poset and let $M$ be a $Q$-module.  Then $M$ is tame if and only if it
		admits (and hence admits all of) the following finite descriptions:
		\begin{enumerate}[label=(\roman*)]
			\item a finite poset encoding subordinate to $M$;
			\item a finite fringe presentation of $M$;
			\item finite resolutions by upset modules and by downset modules.
		\end{enumerate}
		Moreover, the constructions can be chosen compatibly with suitable auxiliary geometric
		structures (e.g.\ semialgebraic or piecewise-linear settings) when $Q$ sits inside a
		partially ordered real vector space, and tame morphisms lift to morphisms of the
		corresponding finite presentations/resolutions.
	\end{theorem}
	
	\begin{remark}[Why equivalences matter for software]
		\label{rem:equivalences-matter}
		Theorem~\ref{thm:syzygy} has two practical interpretations:
		\begin{enumerate}[label=(\roman*)]
			\item It legitimizes switching among data representations depending on what is being
			computed.  One might build a module from geometric data (suggesting a region-based
			encoding), transform it to a fringe presentation to expose ``birth/death'' structure,
			and then compute derived invariants via an indicator resolution.
			\item It explains why it is sensible to base a computational pipeline on finite encodings.
			If a module is tame, then \emph{some} finite encoding exists, and any other finite
			representation can (in principle) be converted to an encoding suitable for computation.
		\end{enumerate}
		TAMER-OP is designed around these equivalences: the library implements multiple entry points
		(fringe-style and region-style) and reduces computations to finite poset models.
	\end{remark}
	
	A more elementary but extremely important consequence is that, once an encoding is fixed,
	many constructions can be carried out on the finite poset $P$ and then transported back
	to $Q$.
	
	\begin{lemma}[Reduction along an encoding]
		\label{lem:reduction-along-encoding}
		Suppose $M\cong \pi^*H$ where $\pi\colon Q\to P$ is a finite encoding and $H$ is a
		$P$-module.
		Then every structure map of $M$ is obtained by applying $H$ to the image under $\pi$:
		for $q\le q'$ in $Q$,
		\[
		M(q\le q') \ = \ H\!\big(\pi(q)\le \pi(q')\big).
		\]
		In particular, any invariant of $M$ that is defined functorially from the diagram of
		vector spaces and structure maps (for example, ranks of structure maps, limits/colimits,
		or derived functors computed from resolutions) can be computed from $H$ on the finite
		poset $P$.
	\end{lemma}
	
	\begin{proof}
		The identity $M=H\circ\pi$ is the definition of pullback.  The final statement follows
		because $\pi^*$ is exact (Lemma~\ref{lem:pullback-exact}) and because the data of $M$
		is completely determined by the data of $H$ together with $\pi$.
	\end{proof}
	
	\subsection{What tameness buys you: algorithmic consequences}
	\label{subsec:tame-algo}
	
	The abstract consequences above translate directly into a computational strategy.
	
	\paragraph{Finite linear algebra.}
	If $P$ is finite and $H$ is a $P$-module with finite-dimensional stalks, then $H$ is
	specified by:
	\begin{enumerate}[label=(\roman*)]
		\item a finite list of vector spaces $H(p)$, and
		\item linear maps $H(p\le p')$ subject only to functoriality constraints.
	\end{enumerate}
	In a finite poset, functoriality can be enforced by recording maps only along cover
	relations (the Hasse diagram) and composing along paths.  This is precisely the kind of
	finite linear-algebraic data structure on which one can implement kernels, images,
	cokernels, ranks, and derived constructions.
	
	\paragraph{Homological algebra in a finite abelian category.}
	The abelian category $\mathbf{Vect}_{\kk}^{P}$ has enough projectives and injectives,
	and its basic building blocks are the principal upset and downset modules
	(Remark~\ref{rem:proj-inj-finite-poset}).  Consequently, one can compute projective and
	injective resolutions by explicit finite algorithms, and then compute $\Tor$ and $\Ext$
	groups by linear algebra on chain complexes.
	When a $Q$-module is tame, Theorem~\ref{thm:syzygy} guarantees that such finite
	resolutions exist \emph{upstairs} in $Q$ as well (in the appropriate tame category),
	but the key computational reduction is that it suffices to resolve the encoded object on
	a finite poset model.
	
	\subsection{Bridge to implementation}
	\label{subsec:tame-bridge}
	
	The remainder of this thesis is concerned with turning the preceding theory into a
	usable computational toolkit.  The next sections explain:
	\begin{enumerate}[label=(\roman*)]
		\item which concrete finite presentations and encodings TAMER-OP supports,
		\item how finite encoding posets are constructed from geometric or algebraic input,
		\item how invariants and derived functors are computed on finite poset models, and
		\item how these computations feed into a statistics pipeline for data analysis.
	\end{enumerate}
	In short: tameness provides the \emph{mathematical contract} that makes the entire
	data-ingestion-to-statistics pipeline finite and computable.
	
	
	
	\section{Overview of TAMER-OP}
	

	
	
	\subsection{Design goals stated mathematically}
	
	
	
	\subsection{The end-to-end pipeline}
	
	\subsection{Core mathematical objects in the API}
	
	\subsection{Architectural layering}
	
	\subsection{Reproducibility and correctness}
	
	\section{The presentation formats}\label{sec:presentation-formats}
	
	The aim of our work is to smoothen a computational tension in the
	mathematics of tame persistence modules.  On the one hand, the natural indexing posets
	that arise in applications are typically infinite (e.g.\ $\mathbb{Z}^n$ or $\mathbb{R}^n$),
	and a persistence module is, definitionally, an infinite functor diagram.  On the other hand,
	any concrete computation must begin from \emph{finite data}.  Miller's syzygy theorem
	(Theorem~\ref{thm:syzygy}) formalizes the central promise of tameness: a tame module admits
	multiple \emph{equivalent} finite descriptions (finite encodings, finite fringe presentations,
	finite indicator resolutions), and one is free to choose whichever description is most
	computationally convenient for a given task \cite{Mil20}.
	
	TAMER-OP turns that promise into a set of concrete data structures and algorithms.  This section describes the
	finite presentation formats implemented in the codebase and explains how they relate to the
	underlying theory.  It is useful to separate three conceptual layers:
	\begin{enumerate}[label=(\roman*)]
		\item \textbf{Front-end (Input formats).}  These are finite objects that arise naturally
		from prior constructions and that encode a tame $Q$-module for an infinite or large poset $Q$.
		\item \textbf{Encodings (compression to a finite poset).}  From a front-end one constructs an
		order-preserving encoding map $\pi\colon Q\to P$ to a finite poset $P$, and pushes the module down
		to a $P$-module. 
		\item \textbf{Finite poset modules (computational core).}  Once the module lives on a finite poset
		$P$, all downstream computations reduce to finite linear algebra (kernels, images, Ext/Tor, rank queries,
		M\"obius inversion, slicing, and so on).  This finite core object is the common target of the conversions
		in this section and the common input to the algorithms in later sections.
	\end{enumerate}
	
	The present section concerns (i), with an emphasis on how each front-end is designed to feed (ii)
	without forcing an early discretization (such as committing to a large bounding box in $\mathbb{Z}^n$
	or an overly fine grid in $\mathbb{R}^n$).  Although we mention the conversion routines that produce
	encodings, the encoding algorithms themselves are deferred to Section~\ref{sec:finite-encoding-constructors}.
	
	\subsection{Philosophy: multiple equivalent front-ends}\label{subsec:formats-philosophy}
	
	A practical library for multiparameter persistence needs to accept input from several sources.
	For example:
	\begin{itemize}
		\item A user might already have a \emph{finite} poset model (an encoding poset built elsewhere) and wants
		to compute invariants on it.
		\item A discrete pipeline (e.g.\ a $\mathbb{Z}^n$ multifiltration) naturally produces multigraded data
		resembling commutative algebra, where presentations and resolutions are most naturally expressed using
		specialized building blocks on $\mathbb{Z}^n$.
		\item A continuous pipeline (e.g.\ a PL or semialgebraic construction over $\mathbb{R}^n$) often yields
		piecewise-constant behavior on polyhedral regions, and it is unnatural (and numerically risky) to snap this
		immediately to a grid.
	\end{itemize}
	
	The unifying observation, following \cite{Mil20}, is that these sources can all be expressed in a common
	pattern: \emph{indicator modules for upward-closed and downward-closed sets, glued by a matrix}.
	This motivates the following heuristic definition.
	
	\begin{defn}[Presentation object in TAMER-OP]\label{def:presentation-object}
		Fix a poset $Q$.  A \emph{presentation object} for a $Q$-module is a finite package of data that
		determines a tame $Q$-module functorially, in the sense that:
		\begin{enumerate}[label=(\roman*)]
			\item for each parameter value $q\in Q$, one can compute, even if implicitly, the fiber $M(q)$,
			and for comparable $q\le q'$ one can compute the structure map $M(q\le q')$;
			\item the module can be encoded by a finite constant subdivision obtained from the same data;
			\item the representation is stable under the basic operations used in the library
			(direct sums, restriction to chains, and the finite encoding step).
		\end{enumerate}
	\end{defn}
	
	This definition is intentionally not a purely mathematical one; it is meant to capture what the code must
	accomplish.  In practice, a presentation object is something like a \texttt{struct} whose fields are (i) a finite
	list of ``birth'' shapes (upsets), (ii) a finite list of ``death'' shapes (downsets), and (iii) a scalar
	matrix specifying how births map into deaths.  The precise meaning of ``shape'' depends on the ambient poset:
	\begin{itemize}
		\item for a finite poset $P$, shapes are subsets of the vertex set, represented efficiently as bitmasks;
		\item for $\mathbb{Z}^n$, shapes are the indecomposable \emph{flats} and \emph{injectives} from Miller's
		$\mathbb{Z}^n$ theory \cite{Mil20};
		\item for $\mathbb{R}^n$ in PL settings, shapes are unions of polyhedral upsets/downsets, represented by
		rational half-space data (or by specialized axis-aligned primitives).
	\end{itemize}
	
	A key design decision, again following Miller's principles, is that TAMER-OP does \emph{not} attempt to compute invariants directly on the infinite posets $\mathbb{Z}^n$ or $\mathbb{R}^n$ from these presentations.  Instead, the role of the presentation is to provide enough structure to build a finite encoding poset $P$ and an encoding map $\pi\colon Q\to P$. Once $P$ is available, the library works with the explicit $P$-modules and standard finite algorithms in the abelian category $\mathbf{Vect}_\kk^P$ (Section~\ref{sec:derived-functors}).
	
	\subsection{Finite fringe presentations}\label{subsec:finite-fringe-presentations}
	
	The most literal computational realization of Miller's framework is obtained when the indexing poset is
	already finite.  In that case, a $P$-module is simply a finite diagram of vector spaces and linear maps, and
	indicator modules for upsets/downsets become finitely supported objects.  This case plays two roles in TAMER-OP:
	\begin{enumerate}[label=(\roman*)]
		\item it provides a clean, fully explicit core format for correctness testing and for downstream algebra; and
		\item it is the \emph{output} format of the encoding step, even when the original module lived on
		$\mathbb{Z}^n$ or $\mathbb{R}^n$.
	\end{enumerate}
	
	\subsubsection*{Finite posets and bitmask upsets/downsets}
	
	In the codebase (see \texttt{FiniteFringe.jl}), a finite poset $P$ is represented by a finite set of vertices
	$\{1,\dots,N\}$ together with a cached order relation.  Upsets and downsets are represented by bit-vectors of
	length $N$.
	
	\begin{defn}[Finite-poset upsets/downsets as data]\label{def:finite-upset-downset-data}
		Let $P$ be a finite poset with vertex set $\{1,\dots,N\}$.
		A \emph{finite upset datum} is a bitmask $U\in\{0,1\}^N$ such that
		$U(p)=1$ and $p\le q$ imply $U(q)=1$.
		A \emph{finite downset datum} is a bitmask $D\in\{0,1\}^N$ such that
		$D(q)=1$ and $p\le q$ imply $D(p)=1$.
	\end{defn}
	
	In addition to explicit bitmasks, the implementation supports closures of arbitrary subsets:
	given a list of generators, one can compute the smallest upset (resp.\ downset) containing them by repeatedly
	closing upward (resp.\ downward).  This is useful when constructing projective/injective covers later, because
	principal upsets $\kk[\uparrow p]$ and principal downsets $\kk[\downarrow p]$ appear repeatedly
	(cf.\ Remark~\ref{rem:proj-inj-finite-poset}).
	
	\subsubsection*{Fringe data and degreewise evaluation}
	
	The main presentation format on a finite poset is a finite fringe presentation.
	
	\begin{defn}[Finite fringe presentation data]\label{def:finite-fringe-data}
		Let $P$ be a finite poset and fix a field $\kk$.
		A \emph{finite fringe presentation datum} consists of:
		\begin{enumerate}[label=(\roman*)]
			\item a finite list of upsets $U_1,\dots,U_m\subseteq P$ (birth regions),
			\item a finite list of downsets $D_1,\dots,D_r\subseteq P$ (death regions),
			\item a matrix $\Phi\in \kk^{r\times m}$ with the \emph{support condition}
			$\Phi_{ji}=0$ whenever $U_i\cap D_j=\varnothing$,
		\end{enumerate}
		and it defines a $P$-module as the image of the induced map
		\[
		\varphi \colon \bigoplus_{i=1}^m \kk[U_i] \longrightarrow \bigoplus_{j=1}^r \kk[D_j]
		\]
		whose component maps $\kk[U_i]\to \kk[D_j]$ are multiplication by $\Phi_{ji}$ on
		$U_i\cap D_j$ and zero elsewhere (compare \cite[Defs.~3.16--3.17]{Mil20}).
	\end{defn}
	
	Degreewise evaluation is particularly direct.  For a vertex $p\in P$, define the active column set
	$I(p)=\{i : p\in U_i\}$ and active row set $J(p)=\{j : p\in D_j\}$.  Evaluating $\varphi$ at $p$ yields a linear
	map
	\[
	\varphi_p \colon \kk^{I(p)} \longrightarrow \kk^{J(p)}
	\]
	whose matrix is the submatrix $\Phi_{J(p),I(p)}$.  The fiber of the module at $p$ is
	\[
	M(p)=\im(\varphi_p)
	\]
	so
	\[
	\dim_\kk M(p)=\mathrm{rank}\big(\Phi_{J(p),I(p)}\big).
	\]
	In TAMER-OP, this is implemented directly: membership in each $U_i$ and $D_j$ is a bit-test, so constructing
	$I(p)$ and $J(p)$ is efficient, and rank computations are performed over exact coefficients (typically $\mathbb{Q}$)
	to avoid numerical rank instability.
	
	\begin{remark}[Why the support condition matters computationally]
		The support condition $\Phi_{ji}=0$ when $U_i\cap D_j=\varnothing$ is the finite-poset analog of a very
		practical constraint: it guarantees that evaluation at each $p$ never introduces spurious maps between
		zero and nonzero stalks.  In the code, constructors enforce this condition and can optionally store $\Phi$
		as a sparse matrix when the support is sparse.  This matters because, even when $m$ and $r$ are moderate,
		the dense matrix $\Phi$ can be the dominant memory cost.
	\end{remark}
	
	\subsubsection*{From fringe presentations to explicit $P$-modules}
	
	A fringe presentation is already a finite representation, but many downstream algorithms prefer an explicit
	representation of a $P$-module in terms of:
	\begin{enumerate}[label=(\roman*)]
		\item dimensions $\dim M(p)$ at each vertex $p$, and
		\item linear maps along cover relations in the Hasse diagram, with all other structure maps obtained by
		composition.
	\end{enumerate}
	This explicit format is implemented in the codebase as the \texttt{PModule} type (defined in \texttt{Modules.jl}).
	The conversion \texttt{pmodule\_from\_fringe} (in \texttt{IndicatorResolutions.jl}) follows the mathematics
	suggested by Proposition~\ref{prop:functor-category-abelian}: choose bases for each fiber $M(p)$ and then define the
	map $M(p\le q)$ as the unique linear map making the relevant inclusion squares commute.
	
	Algorithmically, the conversion has two conceptual steps:
	\begin{enumerate}[label=(\roman*)]
		\item For each vertex $p$, compute a basis for the image $\im(\varphi_p)\subseteq \kk^{J(p)}$.
		In exact arithmetic this can be done by row-reduction of $\Phi_{J(p),I(p)}$ and extracting a column-space basis.
		\item For each cover relation $p\prec q$ in $P$, compute the induced linear map
		$M(p)\to M(q)$ by expressing the image basis vectors at $p$ in the basis at $q$ via the natural inclusions
		$\kk^{J(p)}\hookrightarrow \kk^{J(q)}$ coming from downset functoriality.
	\end{enumerate}
	To keep this step practical, the implementation caches intermediate basis computations and reuses them when multiple
	edges share endpoints.  This conversion is also the place where exact arithmetic over $\mathbb{Q}$ is most valuable:
	small numerical errors in rank computations would otherwise lead to inconsistent dimensions and broken functoriality.
	
	\subsection{$\mathbb{Z}^n$ presentations}\label{subsec:Zn-presentations}
	
	When the indexing poset is $Q=\mathbb{Z}^n$ with coordinatewise order, persistence modules become closely related to
	multigraded commutative algebra (Section~\ref{sec:background:viewpoints}).  Concretely, $\mathbb{N}^n$-indexed modules
	are equivalent to multigraded modules over $\kk[x_1,\dots,x_n]$ (Proposition~\ref{prop:graded-module-equivalence}),
	and even for $\mathbb{Z}^n$-indexing, many of the same structural ideas persist \cite{CZ09,Oud15}.
	
	A naive approach to computation on $\mathbb{Z}^n$ would restrict to a large finite box and treat it as a finite poset.
	This is workable for toy problems but scales poorly, since the number of lattice points grows exponentially in $n$.
	TAMER-OP therefore uses a presentation format tailored to Miller's $\mathbb{Z}^n$ theory: \emph{flange presentations}
	built from indecomposable \emph{flats} and \emph{injectives} \cite[\S5]{Mil20}.  The guiding idea is the same as in
	finite fringe presentations, but the building blocks are chosen so that membership and intersection are simple and so
	that finite encodings can be constructed without enumerating a huge grid.
	
	\subsubsection*{Faces, flats, and injectives}
	
	Fix $n\ge 1$ and write $\mathbf{e}_1,\dots,\mathbf{e}_n$ for the standard basis of $\mathbb{Z}^n$.
	A \emph{face} (in the sense used in the code and in \cite{Mil20}) is a choice of a subset of coordinates along which
	we allow bi-infinite translation.
	
	\begin{defn}[Face subgroup]\label{def:face-subgroup}
		Let $\tau\subseteq \{1,\dots,n\}$.
		Define the subgroup
		\[
		\mathbb{Z}^\tau := \{v\in \mathbb{Z}^n : v_i=0 \text{ for } i\notin \tau\}.
		\]
		Equivalently, $\mathbb{Z}^\tau$ is spanned by $\{\mathbf{e}_i : i\in\tau\}$.
	\end{defn}
	
	\begin{defn}[Indecomposable flat and injective shapes]\label{def:flat-injective-shapes}
		Fix $\tau\subseteq\{1,\dots,n\}$ and a basepoint $\mathbf{b}\in\mathbb{Z}^n$.
		Define subsets of $\mathbb{Z}^n$:
		\[
		F(\mathbf{b},\tau) := \mathbf{b} + \mathbb{N}^n + \mathbb{Z}^\tau,
		\qquad
		E(\mathbf{b},\tau) := \mathbf{b} + \mathbb{Z}^\tau - \mathbb{N}^n.
		\]
		Then $F(\mathbf{b},\tau)$ is an upset in $\mathbb{Z}^n$, and $E(\mathbf{b},\tau)$ is a downset.
		The corresponding indicator modules $\kk[F(\mathbf{b},\tau)]$ and $\kk[E(\mathbf{b},\tau)]$ are the basic
		indecomposable ``flat'' and ``injective'' building blocks used in flange presentations \cite{Mil20}.
	\end{defn}
	
	In the Julia implementation (\texttt{FlangeZn.jl}), a face $\tau$ is stored as a bitmask, and a flat/injective is stored as the pair $(\mathbf{b},\tau)$ together with an optional identifier.  Membership tests reduce to coordinatewise inequalities after projecting away the $\tau$ directions. This makes evaluation cheap even when $n$ is moderate.
	
	\subsubsection*{Flange presentations}
	
	\begin{defn}[Flange presentation over $\mathbb{Z}^n$]\label{def:flange-presentation}
		A \emph{$\mathbb{Z}^n$ flange presentation} (over a field $\kk$) consists of:
		\begin{enumerate}[label=(\roman*)]
			\item a finite list of flat indicators $\kk[F_1],\dots,\kk[F_m]$ with $F_i=F(\mathbf{b}_i,\tau_i)$,
			\item a finite list of injective indicators $\kk[E_1],\dots,\kk[E_r]$ with $E_j=E(\mathbf{c}_j,\sigma_j)$,
			\item a scalar matrix $\Phi\in \kk^{r\times m}$ (rows indexed by injectives, columns by flats),
		\end{enumerate}
		which determines a $\mathbb{Z}^n$-module
		\[
		M := \im\!\left(\varphi \colon \bigoplus_{i=1}^m \kk[F_i] \longrightarrow \bigoplus_{j=1}^r \kk[E_j]\right),
		\]
		where the component map $\kk[F_i]\to \kk[E_j]$ is multiplication by $\Phi_{ji}$ on $F_i\cap E_j$ and zero otherwise
		(compare \cite[Rems.~5.5 and 5.11, Def.~5.14]{Mil20}).
	\end{defn}
	
	The degreewise evaluation mirrors the finite-poset case quite closely.  For $\mathbf{g}\in\mathbb{Z}^n$,
	let $I(\mathbf{g})=\{i : \mathbf{g}\in F_i\}$ and $J(\mathbf{g})=\{j : \mathbf{g}\in E_j\}$.
	Then evaluation at $\mathbf{g}$ yields a linear map
	\[
	\varphi_{\mathbf{g}} \colon \kk^{I(\mathbf{g})} \longrightarrow \kk^{J(\mathbf{g})}
	\]
	with matrix $\Phi_{J(\mathbf{g}),I(\mathbf{g})}$, and hence
	\[
	\dim_\kk M(\mathbf{g}) = \mathrm{rank}\big(\Phi_{J(\mathbf{g}),I(\mathbf{g})}\big).
	\]
	In the code, this is implemented by extracting the active row and column indices via membership tests
	and then computing the rank of the resulting submatrix.  For correctness, the default coefficient field used
	throughout the current implementation is $\mathbb{Q}$, with rank computed by exact row reduction.
	
	\begin{remark}[Automatic enforcement of the monomial support condition]
		As with finite fringe presentations, flange presentations should respect a support condition:
		entries corresponding to empty intersections $F_i\cap E_j=\varnothing$ do not contribute to any degree and are
		safely forced to zero.  The \texttt{Flange} constructor in \texttt{FlangeZn.jl} enforces this by explicitly
		zeroing such entries, which prevents subtle bugs where a user accidentally supplies a dense matrix with
		spurious coefficients.  This is also a mild but useful performance improvement: it increases sparsity in
		common cases where the intersection pattern is already sparse.
	\end{remark}
	
	\subsubsection*{Why flanges are the preferred $\mathbb{Z}^n$ front-end}
	
	Flange presentations are not the only finite representation of tame $\mathbb{Z}^n$-modules, but they have several
	advantages that justify their role as the primary $\mathbb{Z}^n$ front-end:
	\begin{itemize}
		\item \textbf{They are native to the theory.}  They are the $\mathbb{Z}^n$ analog of fringe presentations and arise naturally from Miller's development of finitely determined modules \cite{Mil20}.
		\item \textbf{They avoid bounding boxes.}  One can often build a finite encoding poset using only the finite list
		of flats and injectives, rather than enumerating a large grid of degrees.  This design goal strongly influences
		the encoding algorithms in Section~\ref{sec:finite-encoding-constructors}.
		\item \textbf{They are compatible with commutative algebra tooling.}  In multigraded algebra, flat and injective
		objects play roles analogous to localizations and injective hulls; they align well with resolution-based invariants.
	\end{itemize}
	The cost is that a flange presentation is less immediately geometric than a polyhedral $\mathbb{R}^n$ presentation.
	TAMER-OP therefore treats flanges as a specialized, highly efficient format for lattice-indexed modules, not as a
	universal one.
	
	\subsection{$\mathbb{R}^n$ and piecewise-linear presentations}\label{subsec:Rn-PL-presentations}
	
	Many persistence modules arising from real-parameter filtrations are \emph{piecewise constant} on regions cut out by a finite set of geometric constraints.  In the simplest PL settings, these constraints come from finitely many affine half-spaces and their intersections.  This observation is reflected both in the theory of constructible persistence modules and in sheaf-theoretic approaches \cite{KS18,Mil23}.  In TAMER-OP, the corresponding front-end is a \emph{PL fringe presentation}: a fringe presentation whose upsets and downsets are represented as polyhedral regions in $\mathbb{R}^n$ (or, in a specialized fast path, as unions of axis-aligned boxes).
	
	\subsubsection*{Geometric upsets and downsets}
	
	Let $Q=\mathbb{R}^n$ with the coordinatewise order.
	A subset $U\subseteq \mathbb{R}^n$ defines an upset indicator module $\kk[U]$ only when $U$ is upward closed;
	likewise $D$ must be downward closed.  For computation, we need concrete representations of such sets and efficient
	membership tests.
	
	The codebase supports two PL backends (see \texttt{PLPolyhedra.jl} and \texttt{PLBackend.jl}), reflecting a  tradeoff:
	\begin{itemize}
		\item \textbf{General polyhedral backend.}  Upsets and downsets are represented as finite unions of rational
		convex polyhedra (typically specified by half-space inequalities).  This is flexible and matches the theoretical
		PL picture but can require heavier polyhedral geometry routines.
		\item \textbf{Axis-aligned box backend.}  Upsets and downsets are represented as unions of axis-aligned boxes
		(rectangles in $n=2$).  This case supports very fast region decompositions and is well adapted to many practical
		filtrations where constraints are coordinate thresholds.
	\end{itemize}
	In both backends, the intent is that the sets are \emph{certifiably} upward or downward closed.  For the purposes of
	this thesis, it suffices to emphasize that each shape supports:
	\begin{enumerate}[label=(\roman*)]
		\item membership testing $q\in U_i$ or $q\in D_j$,
		\item intersection testing (needed to enforce the matrix support condition),
		\item extraction of a finite set of ``boundary features'' that will generate the constant subdivision used for
		encoding.
	\end{enumerate}
	The encoding algorithms that exploit these features are postponed to Section~\ref{sec:finite-encoding-constructors}.
	
	\subsubsection*{PL fringe presentations}
	
	With geometric upsets and downsets in hand, the PL fringe presentation is formally identical to the finite fringe
	presentation, except that $Q=\mathbb{R}^n$ is infinite and the shapes are geometric rather than combinatorial.
	
	\begin{defn}[PL fringe presentation]\label{def:pl-fringe-presentation}
		A \emph{PL fringe presentation} over $Q=\mathbb{R}^n$ consists of:
		\begin{enumerate}[label=(\roman*)]
			\item finitely many geometric upsets $U_1,\dots,U_m\subseteq\mathbb{R}^n$,
			\item finitely many geometric downsets $D_1,\dots,D_r\subseteq\mathbb{R}^n$,
			\item a matrix $\Phi\in\kk^{r\times m}$ satisfying $\Phi_{ji}=0$ whenever $U_i\cap D_j=\varnothing$,
		\end{enumerate}
		and it defines the $\mathbb{R}^n$-module
		\[
		M := \im\!\left(\bigoplus_{i=1}^m \kk[U_i] \xrightarrow{\ \Phi\ } \bigoplus_{j=1}^r \kk[D_j]\right).
		\]
	\end{defn}
	
	As before, the fiber dimension at $q\in\mathbb{R}^n$ is computed by restricting $\Phi$ to active rows and columns and
	taking a rank.  In practice, however, such pointwise queries are only a diagnostic tool.  The primary use of a PL
	fringe presentation is as input to a finite encoding step: the shapes $U_i$ and $D_j$ determine a finite subdivision of
	$\mathbb{R}^n$ into regions on which all membership patterns are constant, and hence on which $M$ is constant.  This
	recovers the ``finite regions of behavior'' intuition behind tameness \cite{Mil20}.
	
	\begin{remark}[Exactness and geometric robustness]
		The PL setting is where numerical issues most easily corrupt a computation, because region boundaries are defined by equalities and inequalities.  The current implementation therefore emphasizes rational geometry where possible: polyhedral constraints are stored with rational coefficients, and rank computations are performed over $\mathbb{Q}$.  This does not eliminate all geometric subtleties, as degenerate intersections still require care. However, it does avoid potential failures in persistent computations due to incorrect ranks induced by floating-point instability near boundary cases.
	\end{remark}
	
	\subsection{Conversions between formats}\label{subsec:format-conversions}
	
	The core promise of this section is that the front-ends described above are not siloed.
	They are designed so that, after a finite encoding is constructed, they can all be converted to the same finite core
	format.  The overarching conversion pattern is:
	
	\begin{center}
		\begin{tabular}{c}
			(front-end presentation on $Q$)\\
			$\Downarrow$ (construct finite encoding $\pi\colon Q\to P$)\\
			(fringe presentation on finite $P$)\\
			$\Downarrow$ (basis extraction + cover-edge maps)\\
			(explicit $P$-module \texttt{PModule})
		\end{tabular}
	\end{center}
	
	The theoretical justification is that the encoding poset captures the module up to isomorphism in the tame category
	(Definition~\ref{def:tame-morphism} and Theorem~\ref{thm:tame-abelian}), and all further invariants computed by TAMER-OP
	factor through this finite model (Lemma~\ref{lem:reduction-along-encoding}).
	
	\subsubsection*{Finite-poset fringe $\to$ explicit $P$-module}
	
	This is the simplest conversion and was already outlined in Subsection~\ref{subsec:finite-fringe-presentations}.
	It is implemented by \texttt{pmodule\_from\_fringe}, which computes image bases at each vertex and induced maps on cover
	relations.
	
	\subsubsection*{$\mathbb{Z}^n$ flange $\to$ finite encoding $\to$ $P$-module}
	
	For $\mathbb{Z}^n$ flanges, the key step is constructing an encoding without enumerating a large portion of the lattice.
	At a conceptual level, the conversion uses the finite family of sets
	\[
	\{F_i\}_{i=1}^m \ \cup\ \{\mathbb{Z}^n\setminus E_j\}_{j=1}^r
	\]
	to define a \emph{signature} map
	\[
	\mathrm{sig}\colon \mathbb{Z}^n \longrightarrow \{0,1\}^{m+r},
	\qquad
	\mathrm{sig}(\mathbf{g}) = \big( \mathbf{1}_{\mathbf{g}\in F_1},\dots,\mathbf{1}_{\mathbf{g}\in F_m},
	\mathbf{1}_{\mathbf{g}\notin E_1},\dots,\mathbf{1}_{\mathbf{g}\notin E_r}\big).
	\]
	Points with the same signature lie in the same constant region for the presentation, hence the same fiber of an
	encoding map.  The nontrivial work is to show that only finitely many signatures occur under suitable finiteness
	hypotheses and to construct a finite poset $P$ of regions ordered by reachability.  This is precisely the kind of
	finite constant subdivision that Miller uses to define tame morphisms and encodings \cite{Mil20}.  The algorithmic
	realization in the codebase is implemented in \texttt{ZnEncoding.jl}.
	
	Once $P$ and $\pi$ are constructed, the flange is pushed down to a fringe presentation on $P$ by evaluating each flat
	and injective region on a representative for each region of $P$, producing explicit upsets and downsets in the finite
	poset.  The remainder of the pipeline (fringe $\to$ \texttt{PModule}) is then identical to the finite case.
	
	\subsubsection*{PL fringe $\to$ finite encoding $\to$ $P$-module}
	
	The PL case follows the same signature philosophy, but now using geometric objects.  The constant subdivision is obtained from the finite set of boundaries that define the polyhedral upsets and downsets, yielding finitely many regions when the input is PL or semialgebraic.  The polyhedral backend (\texttt{PLPolyhedra.jl}) and the axis-aligned backend (\texttt{PLBackend.jl}) implement different concrete strategies for producing the same output: a finite poset $P$, an encoding map $\pi\colon \mathbb{R}^n\to P$ (with a point-location routine), and a pushed-down fringe presentation on $P$.
	
	\begin{remark}[Why conversions are staged]
		It is tempting to attempt a direct conversion from an infinite presentation (flange or PL fringe) to an explicit
		$P$-module by sampling many points and inferring structure maps.  We avoid this on TAMER-OP.  The staged approach (front-end $\to$ encoding $\to$ fringe on $P$ $\to$ explicit $P$-module) mirrors the theory and isolates sources of error. Here, any mismatch can be diagnosed at the level of signatures and regions before committing to linear algebra on the finite poset.
	\end{remark}
	
	\subsection{Validation}\label{subsec:format-validation}
	
	Compact presentations are powerful precisely because they hide large implicit structure. But, that also makes them fragile. A small inconsistency in the input can produce a module that is not functorial, not tame, or simply not the module the user intended.  TAMER-OP therefore incorporates validation at multiple layers.  The checks fall into two categories: \emph{structural} checks that are inexpensive and always enabled in constructors, and \emph{diagnostic} checks that can be toggled because they may be expensive.
	
	\subsubsection*{Structural checks}
	
	\begin{itemize}
		\item \textbf{Poset correctness.}  For a finite poset, we check that the stored relation is reflexive and transitive and that it is antisymmetric up to equality of vertices.  Cover relations are derived from the order and are used as the canonical generating set of morphisms.
		\item \textbf{Upset/downset closure.}  For finite-poset upsets and downsets represented as bitmasks, constructors
		and closure routines ensure upward-closed and downward-closed properties by construction.
		\item \textbf{Monomial (support) conditions.}  For each presentation format, nonzero matrix entries are forbidden
		when the corresponding birth and death shapes do not intersect.  Constructors enforce this condition
		automatically:
		in the finite fringe format, a check verifies $\Phi_{ji}=0$ when $U_i\cap D_j=\varnothing$; in the $\mathbb{Z}^n$
		flange format, the constructor explicitly zeros such entries; and in the PL formats, intersection tests serve the
		same role.  These checks directly implement the support constraints that make a matrix correspond to an actual
		natural transformation of indicator modules \cite[Prop.~3.18]{Mil20}.
	\end{itemize}
	
	\subsubsection*{Diagnostic checks (optional)}
	
	\begin{itemize}
		\item \textbf{Spot-check functoriality.}  For presentations that support evaluation, one can sample comparable
		pairs $q\le q'$ and verify that the induced maps satisfy basic functorial constraints along short chains.
		\item \textbf{Encoding correctness tests.}  After building an encoding $\pi\colon Q\to P$, we can sample points
		$q,q'\in Q$ with $\pi(q)=\pi(q')$ and verify that the presentation evaluates to the same fiber data at $q$ and $q'$.
		This corresponds directly to the mathematical requirement that $M$ be constant on fibers of $\pi$
		(Remark~\ref{rem:encoding-fibers}).
		\item \textbf{Rank certification.}  When working over $\mathbb{Q}$, ranks are computed by exact elimination; this
		can be used to certify that computed fiber dimensions and structure maps are consistent.  In later sections, the
		same philosophy is used to certify minimality of resolutions on finite posets.
	\end{itemize}
	
	\subsubsection*{Bridge to the encoding algorithms}
	
	At this point we have described the finite data structures that represent tame modules before encoding.
	The next section addresses the algorithmic core of TAMER-OP: given one of these presentation formats, how does one
	construct a finite encoding poset $P$ and an encoding map $\pi\colon Q\to P$ efficiently and robustly, and how does one
	push the module down to $P$ so that all subsequent invariants can be computed by finite linear algebra?
	
	% ============================================================
	% Section 6: Finite encoding constructors
	% ============================================================
	
	\section{Finite encoding constructors}
	\label{sec:finite-encoding-constructors}
	
	A recurring theme in Chapters~\ref{sec:background}--\ref{sec:tame-modules} is that ``tameness'' is not merely a finiteness hypothesis but a promise that persistence modules admit \emph{finite combinatorial models}.  In Miller's language, a tame $Q$-module $M$ admits a \emph{finite poset encoding} $\pi\colon Q\to P$ with
	$P$ finite and $M\cong \pi^*H$ for some $P$-module $H$ (Definition~\ref{def:poset-encoding}), and this is equivalent to the existence of finite fringe-style descriptions (Theorem~\ref{thm:syzygy}) \cite{Mil20}. The purpose of this section is to explain how TAMER-OP \emph{constructs} such a finite encoding from the finite presentations introduced in the previous section.
	
	A fringe or flange presentation is a finite collection of geometric regions in parameter space (upsets and downsets, often polyhedral) together with a matrix of coefficients.  The encoding problem asks for a finite poset $P$ of ``regions of constant behavior'' together with a map $\pi\colon Q\to P$ that collapses each constant region to a single element. Once this is achieved, the computational task shifts decisively: all subsequent operations on $M$ (kernels, images, ranks, resolutions, derived functors, and so on) reduce to finite linear algebra on the encoded $P$-module, as in Lemma~\ref{lem:reduction-along-encoding}.
	
	Algorithmically, TAMER-OP implements encodings in three ambient regimes:
	\begin{itemize}
		\item \textbf{Lattice encodings} for $Q=\mathbb{Z}^n$ (implemented in
		\texttt{ZnEncoding.jl}), where presentations are built from indicator flats and
		injectives (Section~\ref{subsec:Zn-presentations}).
		\item \textbf{Axis-aligned piecewise-linear encodings} for $Q=\mathbb{R}^n$
		with coordinatewise order and regions described by finite unions of orthants
		(implemented in \texttt{PLBackend.jl}).
		\item \textbf{General polyhedral PL encodings} for $Q=\mathbb{R}^n$ where
		regions are unions of rational polyhedra (implemented in \texttt{PLPolyhedra.jl}),
		used when axis-aligned structure is insufficient.
	\end{itemize}
	All three share a single organizing principle: \emph{compute a finite partition of
		parameter space on which the relevant indicator data are constant, then collapse
		cells with the same signature.}
	
	\subsection{The encoding problem}
	\label{subsec:encoding-problem}
	
	Fix an ambient poset $Q$ (typically $Q=\mathbb{Z}^n$ or $Q=\mathbb{R}^n$ with
	coordinatewise order).  The immediate input to an encoding constructor in
	TAMER-OP is a finite presentation in one of the formats from
	Section~\ref{sec:presentation-formats}.  For the present section it is convenient
	to phrase this input uniformly as follows.
	
	\paragraph{Input data model.}
	A finite fringe-style presentation provides:
	\begin{enumerate}[label=(\roman*)]
		\item a finite collection of \emph{birth upsets} $U_1,\dots,U_m\subseteq Q$,
		\item a finite collection of \emph{death downsets} $D_1,\dots,D_r\subseteq Q$,
		\item a matrix $\Phi\in \kk^{r\times m}$ specifying a map
		$\varphi\colon \bigoplus_{i=1}^m \kk[U_i]\to \bigoplus_{j=1}^r \kk[D_j]$
		(Lemmas~\ref{lem:monomial-support} and Definition~\ref{def:monomial-matrix}).
	\end{enumerate}
	In the lattice implementation these $U_i$ and $D_j$ are not literally stored as
	arbitrary subsets; instead they come from structured families (indicator flats and
	injectives) that admit fast membership tests and finite boundary descriptions
	(\texttt{IndFlat} and \texttt{IndInj} in \texttt{FlangeZn.jl}).  In the PL
	implementations the upsets and downsets are represented as finite unions of
	structured polyhedral sets (orthants in \texttt{PLBackend.jl}, general rational
	polyhedra in \texttt{PLPolyhedra.jl}).
	
	\paragraph{Desired output.}
	An encoding constructor produces:
	\begin{enumerate}[label=(\roman*)]
		\item a finite poset $P$ (the \emph{encoding poset}),
		\item a poset morphism $\pi\colon Q\to P$ (the \emph{encoding map}),
		\item a $P$-module $H\colon P\to \mathbf{Vect}_{\kk}$, such that the presented
		$Q$-module factors as $M\cong \pi^*H$.
	\end{enumerate}
	In code, the output is typically packaged as \texttt{(P, H, pi)} or
	\texttt{(P, Hs, pi)} when encoding several modules at once.
	Here $H$ is stored as a finite poset module in \texttt{FiniteFringe.jl} (usually a
	\texttt{FringeModule} on $P$), and $\pi$ is stored as a lightweight object
	supporting point location queries (for example \texttt{ZnEncodingMap}).
	
	\paragraph{A methodological choice.}
	There is no unique ``best'' encoding in general.  Even when a tame encoding exists,
	it need not be canonical, and the problem of finding a minimal encoding can be
	computationally expensive or ill-posed under different equivalence notions.
	TAMER-OP therefore adopts a \emph{presentation-relative} encoding strategy:
	it constructs an encoding that is guaranteed to be subordinate to the \emph{specific}
	finite family of upsets and downsets appearing in the chosen input format.  This
	choice matches the computational contract of Miller's equivalences:
	theorem-level existence is supplied by tameness \cite{Mil20}, while the encoding
	constructor is a deterministic algorithm whose complexity can be analyzed in terms
	of the size of the given finite presentation.
	
	\subsection{Region/signature idea}
	\label{subsec:signature-idea}
	
	The core encoding mechanism used throughout TAMER-OP can be stated abstractly
	without reference to a particular ambient domain.
	
	\begin{defn}[Signature map associated to a finite fringe interface]
		\label{def:signature-map}
		Let $Q$ be a poset and suppose we are given upsets $U_1,\dots,U_m\subseteq Q$
		and downsets $D_1,\dots,D_r\subseteq Q$.
		Define functions
		\[
		\chi_{U_i}(q) :=
		\begin{cases}
			1 & q\in U_i,\\
			0 & q\notin U_i,
		\end{cases}
		\qquad
		\chi_{\overline{D_j}}(q) :=
		\begin{cases}
			1 & q\notin D_j,\\
			0 & q\in D_j,
		\end{cases}
		\]
		and define the \emph{signature} of $q\in Q$ to be the bit vector
		\[
		\sigma(q) := \big(\chi_{U_1}(q),\dots,\chi_{U_m}(q)\,;\,
		\chi_{\overline{D_1}}(q),\dots,\chi_{\overline{D_r}}(q)\big)\in \{0,1\}^{m+r}.
		\]
	\end{defn}
	
	The reason for encoding downsets via their complements is monotonicity: if $D$ is
	a downset, then $Q\setminus D$ is an upset, hence order-preserving in the same
	direction as birth regions.
	
	\begin{lemma}[Monotonicity of signatures]
		\label{lem:signature-monotone}
		The signature map $\sigma\colon Q\to \{0,1\}^{m+r}$ in
		Definition~\ref{def:signature-map} is order-preserving when
		$\{0,1\}^{m+r}$ is ordered coordinatewise.
	\end{lemma}
	
	\begin{proof}
		If $q\le q'$ in $Q$ and $q\in U_i$ for some upset $U_i$, then $q'\in U_i$ by upward
		closure, so $\chi_{U_i}(q)\le \chi_{U_i}(q')$.
		If $q\le q'$ and $q\notin D_j$ for a downset $D_j$, then $q'\notin D_j$ as well,
		because membership in a downset can only be lost when moving upward.  Equivalently,
		$Q\setminus D_j$ is an upset.  Hence $\chi_{\overline{D_j}}(q)\le
		\chi_{\overline{D_j}}(q')$ for all $j$.
	\end{proof}
	
	The following encoding construction follows naturally from the lemma.
	
	\begin{proposition}[Encoding via realized signatures]
		\label{prop:encoding-by-signatures}
		Let $P:=\mathrm{im}(\sigma)\subseteq \{0,1\}^{m+r}$ with the induced partial order,
		and let $\pi:=\sigma\colon Q\to P$ (followed by the projection onto the image).
		Then for each $i$ and $j$, the indicator modules $\kk[U_i]$ and $\kk[D_j]$ factor
		through $\pi$.  Consequently, any fringe presentation map
		$\varphi\colon \bigoplus_i \kk[U_i]\to \bigoplus_j \kk[D_j]$ induces a
		$P$-module presentation whose pullback along $\pi$ recovers the original one.
	\end{proposition}
	
	\begin{proof}[Proof sketch]
		For each $i$, define an upset $\widehat{U}_i\subseteq P$ by
		\[
		\widehat{U}_i:=\{p\in P : \text{the $i$th upset-bit of $p$ equals }1\}.
		\]
		Because the $i$th coordinate is monotone in $P$, this is indeed an upset in $P$,
		and by construction $\pi^{-1}(\widehat{U}_i)=U_i$.
		Thus $\kk[U_i]\cong \pi^*\kk[\widehat{U}_i]$.
		Similarly, define a downset $\widehat{D}_j\subseteq P$ by declaring $p\in\widehat{D}_j$
		if the $(m+j)$th bit of $p$ (the complement-of-downset bit) equals $0$, so that
		$\pi^{-1}(\widehat{D}_j)=D_j$.
		The map $\varphi$ is specified stalkwise by the matrix $\Phi$ on intersections
		$U_i\cap D_j$, and those intersections are preserved under pullback because each
		$U_i$ and $D_j$ is a union of fibers of $\pi$.
	\end{proof}
	
	Proposition~\ref{prop:encoding-by-signatures} conceptually encompasses encoding as a whole: \emph{once} the set of realized signatures is finite, one has a finite encoding.  The real work is therefore to compute $\mathrm{im}(\sigma)$ efficiently from a structured description of the $U_i$ and $D_j$, and to equip the realized signatures with an implementation-friendly poset structure and a fast point-location map $\pi$.
	
	From a geometric standpoint, this signature construction is a concrete incarnation of the ``constant subdivision'' idea in Miller's definition of tame morphisms \cite{Mil20}.  The regions on which all indicator data are constant form a finite stratification of parameter space, analogous in spirit to the stratifications appearing in constructible sheaf approaches \cite{KS18,Mil23}, though TAMER-OP works directly with the combinatorics of indicator regions rather than with microsupport.
	
	\subsection{Encoding algorithms by ambient domain}
	\label{subsec:encoding-algorithms-by-domain}
	
	This subsection describes how TAMER-OP computes the realized signatures
	$\mathrm{im}(\sigma)$ in each supported regime.  The three backends share the same
	high-level structure but differ in how they represent regions and how they
	enumerate constant cells.
	
	\subsubsection*{(A) Lattice encoding for $Q=\mathbb{Z}^n$}
	\label{subsubsec:encoding-Zn}
	
	The lattice backend is designed for presentations arising from
	$\mathbb{Z}^n$-indexed modules, where indicator regions are built from algebraic
	data that are intrinsically discrete.
	
	\paragraph{Indicator families and membership tests.}
	In \texttt{FlangeZn.jl}, the basic structured regions are:
	\begin{itemize}
		\item \texttt{IndFlat} objects, which encode certain coordinatewise upsets by
		inequalities of the form $g_i\ge b_i$ on a specified subset of coordinates; and
		\item \texttt{IndInj} objects, which encode coordinatewise downsets by inequalities
		$g_i\le b_i$ on a specified subset of coordinates.
	\end{itemize}
	Both admit $O(n)$ membership checks (\texttt{in\_flat} and \texttt{in\_inj}),
	but the key point is that their \emph{boundaries} occur only at finitely many
	integer coordinate values extracted from the $b_i$ in the presentation.
	
	\paragraph{Critical coordinates and a product cell complex.}
	The encoding constructor \texttt{encode\_poset\_from\_flanges} in
	\texttt{ZnEncoding.jl} begins by collecting, for each axis $k=1,\dots,n$, a
	sorted list of \emph{critical coordinates} at which some membership predicate
	changes.  Concretely:
	\begin{itemize}
		\item for each flat constraint $g_k\ge b_k$, the threshold $b_k$ is critical;
		\item for each injective constraint $g_k\le b_k$, the complement threshold
		$b_k+1$ is critical, because in $\mathbb{Z}$ the complement of $g_k\le b_k$
		is $g_k\ge b_k+1$.
	\end{itemize}
	Call these lists $C_k\subset \mathbb{Z}$.  They define a finite product cell
	structure on $\mathbb{Z}^n$ by decomposing each coordinate axis into slabs between
	successive critical values.  The resulting \emph{grid of cells} has size
	\[
	N_{\mathrm{cell}} = \prod_{k=1}^n (|C_k|+1).
	\]
	On each cell, the truth values of all inequalities defining all flats and injectives
	are constant, hence the signature $\sigma$ is constant.
	
	\paragraph{Signature enumeration by incremental scanning.}
	A direct approach would evaluate all $m+r$ membership predicates on a representative
	lattice point of each cell, costing $O((m+r)\,n\,N_{\mathrm{cell}})$ in the worst
	case.  The implementation instead performs a single scan through the grid, updating
	the signature incrementally as it crosses critical coordinates.  The key mechanism
	is:
	\begin{enumerate}[label=(\roman*)]
		\item precompute, for each axis and each critical coordinate, which flats or injective
		complements toggle their constraint status when crossing that coordinate; and
		\item traverse the grid in a snake-like boustrophedon order so that successive cells differ
		in only one coordinate direction.
	\end{enumerate}
	For each cell visited, the scanner updates a small set of counters (how many
	constraints of each flat are still unsatisfied, and how many constraints of each
	injective have been violated) and maintains bitvectors encoding the current
	$(y,z)$ signature.  Each visited cell contributes its signature to a hash table of
	seen signatures, and the scan optionally records a lookup table
	\texttt{cell\_to\_region} mapping cell indices to signature-region indices.
	
	\paragraph{Output.}
	The scan returns the list of realized signatures and a map
	$\pi\colon \mathbb{Z}^n\to P$ implemented as a \texttt{ZnEncodingMap}.  The encoding
	poset $P$ can be represented either:
	\begin{itemize}
		\item implicitly, as a \texttt{SignaturePoset} where $p\le p'$ is tested by subset
		operations on the stored signature bitvectors; or
		\item explicitly, as a \texttt{FinitePoset} whose order relation is materialized as a
		BitMatrix for faster downstream closure operations.
	\end{itemize}
	This choice is controlled by \texttt{poset\_kind} in \texttt{EncodingOptions}.
	
	\subsubsection*{(B) Axis-aligned PL encoding for $Q=\mathbb{R}^n$}
	\label{subsubsec:encoding-PL-boxes}
	
	The axis-aligned PL backend is the continuous analogue of the lattice backend.
	Here $Q=\mathbb{R}^n$ with the coordinatewise order, and the input regions are
	finite unions of \emph{orthants} described by threshold vectors.
	
	\paragraph{Region representation.}
	In \texttt{PLBackend.jl}, an upset is represented as a finite set of ``corners''
	$\ell^{(1)},\dots,\ell^{(s)}\in \mathbb{R}^n$, encoding the union
	\[
	U=\bigcup_{t=1}^s \{x\in \mathbb{R}^n : x\ge \ell^{(t)}\},
	\]
	and similarly a downset is represented by finitely many ``upper corners''
	$u^{(1)},\dots,u^{(t)}$ encoding
	\[
	D=\bigcup_{t=1}^T \{x\in \mathbb{R}^n : x\le u^{(t)}\}.
	\]
	This matches a geometry that should be common in PL persistence constructions, where births and deaths
	are controlled by finitely many coordinate thresholds.
	
	\paragraph{Coordinate hyperplane subdivision.}
	The encoding constructor \texttt{encode\_fringe\_boxes} collects, for each axis
	$k$, all threshold coordinates appearing in any $\ell^{(t)}_k$ or $u^{(t)}_k$.
	These values define a subdivision of $\mathbb{R}$ into open intervals, and hence
	a subdivision of $\mathbb{R}^n$ into open axis-aligned boxes (cells).  As in the
	lattice case, the signature $\sigma$ is constant on each cell.
	
	\paragraph{Representative points and strictness.}
	A subtlety in the continuous setting is boundary handling: if a representative point
	lies on a threshold hyperplane, membership in a closed upset/downset can be ambiguous
	relative to the intended open cell.  TAMER-OP resolves this by choosing a point in the
	\emph{interior} of each cell, implemented via a small offset parameter
	\texttt{strict\_eps} in \texttt{EncodingOptions}.  This makes the signature assignment
	robust and deterministic.
	
	\paragraph{Scanning and output.}
	The axis-aligned backend performs a scan through the grid of cells and records
	realized signatures, using essentially the same boustrophedon mechanism as the lattice
	backend (implemented in \texttt{\_collect\_cell\_signatures\_scan}).
	The output is a finite poset $P$ on realized signatures and an encoding map
	$\pi\colon \mathbb{R}^n\to P$ implemented as a \texttt{BoxEncodingMap}.
	
	\subsubsection*{(C) General polyhedral PL encoding}
	\label{subsubsec:encoding-PL-polyhedra}
	
	Axis-aligned encodings are fast and often sufficient, but they do not cover
	general polyhedral regions (for example, regions defined by non-axis-aligned linear
	inequalities).  The polyhedral backend in \texttt{PLPolyhedra.jl} supports upsets and
	downsets given as finite unions of rational polyhedra.
	
	\paragraph{Polyhedral feasibility viewpoint.}
	In principle, one could build a full arrangement of all supporting hyperplanes and
	enumerate its cells, but this can be expensive.  The implementation instead follows
	a feasibility-driven approach: it enumerates candidate signatures and checks whether
	the region defined by that signature is nonempty.
	
	Concretely, for each candidate signature $(y,z)\in\{0,1\}^{m+r}$, one forms the set
	\[
	R(y,z) :=
	\Big(\bigcap_{i:y_i=1} U_i\Big)
	\cap
	\Big(\bigcap_{i:y_i=0} (Q\setminus U_i)\Big)
	\cap
	\Big(\bigcap_{j:z_j=0} D_j\Big)
	\cap
	\Big(\bigcap_{j:z_j=1} (Q\setminus D_j)\Big),
	\]
	and asks whether $R(y,z)$ is nonempty.  Because the $U_i$ and $D_j$ are unions of
	polyhedra, each complement is also a union of polyhedra and halfspace-defined pieces,
	so this feasibility test branches over choices of pieces.  Branches that yield an empty
	intersection are discarded.
	
	\paragraph{Exact arithmetic for correctness.}
	Feasibility and emptiness checks are performed using rational polyhedral computations
	(via \texttt{Polyhedra.jl} and \texttt{CDDLib.jl} in the implementation).
	This is aligned with the ``exactness versus speed'' discussion in
	Section~\ref{sec:background:conventions}: for encoding construction, correctness of
	region adjacency and signature assignment is structurally important, so exact arithmetic
	is often the right default.
	
	\paragraph{Output and point-location implications.}
	The polyhedral backend still outputs a signature poset $P$ (or an explicit finite
	poset) together with an encoding map $\pi$.  The difference is that $\pi$ now needs a
	polyhedral point-location routine.  TAMER-OP stores each realized region together with
	a witness point and uses a two-stage membership test (fast floating-point rejection
	followed by exact verification near boundaries) to implement $\pi$ efficiently; see
	Subsection~\ref{subsec:point-location}.
	
	\subsection{Constructing the finite poset structure}
	\label{subsec:constructing-poset}
	
	Once a finite set of realized signatures has been computed, TAMER-OP must turn it
	into a finite poset $P$ and transport the input presentation to a $P$-module.
	
	\subsubsection*{Signature poset versus explicit finite poset}
	
	Let $\{(y^{(a)},z^{(a)})\}_{a=1}^R$ be the realized signatures, where
	$y^{(a)}\in\{0,1\}^m$ and $z^{(a)}\in\{0,1\}^r$.  The canonical partial order used in
	all backends is the coordinatewise order:
	\[
	a \le_P b
	\quad\Longleftrightarrow\quad
	y^{(a)} \le y^{(b)}\ \text{and}\ z^{(a)} \le z^{(b)}
	\quad\Longleftrightarrow\quad
	\mathrm{supp}(y^{(a)})\subseteq \mathrm{supp}(y^{(b)}),\ 
	\mathrm{supp}(z^{(a)})\subseteq \mathrm{supp}(z^{(b)}).
	\]
	This is exactly the order inherited from $\{0,1\}^{m+r}$.
	
	In code this is realized in two interchangeable forms:
	\begin{itemize}
		\item \textbf{\texttt{SignaturePoset}} (used when \texttt{poset\_kind = :signature}):
		store the bitvectors and implement $\le_P$ by subset checks on bitvectors.
		\item \textbf{\texttt{FinitePoset}} (used when \texttt{poset\_kind = :dense}):
		materialize the full reflexive transitive order relation as an $R\times R$ BitMatrix.
	\end{itemize}
	The second form is larger in memory but can accelerate repeated closure computations
	(upsets, downsets, cover edges), which appear frequently in later invariant
	computation.
	
	\subsubsection*{Transporting indicator data to $P$}
	
	The transport step is the computational instantiation of
	Proposition~\ref{prop:encoding-by-signatures}.  Each birth upset $U_i$ induces an
	upset $\widehat{U}_i\subseteq P$ by reading off the $i$th birth bit in $y$:
	\[
	\widehat{U}_i := \{a\in P : y^{(a)}_i = 1\}.
	\]
	Similarly, each death downset $D_j$ induces a downset $\widehat{D}_j\subseteq P$ by
	\[
	\widehat{D}_j := \{a\in P : z^{(a)}_j = 0\}.
	\]
	(Equivalently, membership in the complement $Q\setminus D_j$ is recorded by $z_j=1$.)
	
	These $\widehat{U}_i$ and $\widehat{D}_j$ are stored as bitsets on $P$ in the
	\texttt{FiniteFringe.FringeModule} representation.
	
	\subsubsection*{Transporting the monomial matrix}
	
	Given $\Phi\in \kk^{r\times m}$ describing the map
	$\bigoplus_i \kk[U_i]\to \bigoplus_j \kk[D_j]$, we must produce a matrix
	$\widehat{\Phi}$ describing
	$\bigoplus_i \kk[\widehat{U}_i]\to \bigoplus_j \kk[\widehat{D}_j]$.
	By Lemma~\ref{lem:monomial-support}, the only structural constraint is:
	\[
	\widehat{\Phi}_{ji} = 0 \quad\text{whenever}\quad \widehat{U}_i\cap \widehat{D}_j=\varnothing.
	\]
	In a signature encoding this intersection test is purely combinatorial: the
	intersection is nonempty if and only if there exists some region $a\in P$ with
	$y^{(a)}_i=1$ and $z^{(a)}_j=0$.  TAMER-OP therefore constructs $\widehat{\Phi}$ by
	retaining the original coefficient $\Phi_{ji}$ when the intersection is nonempty
	and forcing it to zero otherwise.
	
	At a conceptual level, this step is where the encoding ``forgets geometry'' and
	retains only the finite combinatorics needed for algebraic computation.
	
	\subsection{Point location and evaluation map}
	\label{subsec:point-location}
	
	The encoding map $\pi\colon Q\to P$ must be computable on demand, because later
	pipeline stages repeatedly query fibers $H(\pi(q))$ and structure maps along
	comparabilities in $P$.  TAMER-OP therefore treats point location as part of the
	encoding object, not as an afterthought.
	
	\subsubsection*{Grid-based point location (lattice and axis-aligned PL)}
	
	In the lattice and axis-aligned PL backends, $Q$ is subdivided using coordinate
	thresholds, so locating a point reduces to finding which slab index it occupies
	along each axis.
	
	\paragraph{Lattice case.}
	A \texttt{ZnEncodingMap} stores:
	\begin{itemize}
		\item the critical coordinate lists $C_k$ for each axis,
		\item the realized signature dictionary \texttt{sig\_to\_region},
		\item optionally a lookup table \texttt{cell\_to\_region} when the grid is small
		enough to store (controlled by \texttt{max\_lookup\_cells}).
	\end{itemize}
	Given $g\in \mathbb{Z}^n$, one computes the cell index by $n$ binary searches
	(\texttt{searchsortedlast}) in the sorted coordinate lists, costing
	$O\!\left(\sum_{k=1}^n \log |C_k|\right)$ comparisons.
	If \texttt{cell\_to\_region} is present, region lookup is $O(1)$ afterward.
	If it is absent, the implementation recomputes the signature key by evaluating the
	structured flat/injective predicates (still finite and typically small), and then
	looks up the region in a dictionary.  This makes the time per query essentially
	either:
	\[
	O\!\left(n\log K\right)\quad\text{(with lookup table)},\qquad
	O\!\left(n(m+r)\right)\quad\text{(recompute signature)},
	\]
	where $K=\max_k |C_k|$.
	
	\paragraph{Axis-aligned PL case.}
	A \texttt{BoxEncodingMap} works analogously with real thresholds.  Point location
	again uses $n$ binary searches in sorted coordinate lists, followed by a cell-to-region
	lookup.
	
	\subsubsection*{Polyhedral point location}
	
	In the polyhedral backend, the region fibers are not axis-aligned boxes but rational
	polyhedra (or finite unions thereof).  A naive point location routine would test
	membership against every region, which is worst-case $O(R)$ region tests per query,
	with each region test involving $O(Hn)$ inequality evaluations for $H$ halfspaces.
	TAMER-OP implements a more structured routine in \texttt{PLPolyhedra.jl}:
	\begin{itemize}
		\item each realized region stores a witness point and a floating-point approximation
		of its inequalities;
		\item queries are first filtered using cheap floating-point inequality checks that
		rapidly reject most regions;
		\item candidates that survive the float filter are verified using exact rational
		arithmetic near boundaries.
	\end{itemize}
	This does not change the worst-case asymptotics but substantially improves expected
	performance in regimes where regions are moderately separated or when a coarse
	candidate set can be obtained from coordinate heuristics.
	
	\subsection{Building a common encoding for multiple modules}
	\label{subsec:common-encoding}
	
	Most downstream tasks in TAMER-OP are comparative: one wants to compute invariants
	for a \emph{collection} of modules and then feed those outputs into statistical
	summaries.  For this it is highly advantageous, and often necessary, that all modules
	be encoded over a \emph{common} finite poset $P$ so that their invariants are computed
	on a shared combinatorial domain.
	
	\paragraph{Common encoding by union of generator data.}
	Given finitely many modules $M^{(1)},\dots,M^{(T)}$ presented by upsets/downsets
	$\{U_i^{(t)}\}$ and $\{D_j^{(t)}\}$, TAMER-OP constructs a common encoding by forming
	the combined signature map associated to the union of all these regions.  In practice,
	this means:
	\begin{enumerate}[label=(\roman*)]
		\item concatenate the birth and death region lists across all modules,
		\item run the chosen backend encoding algorithm once on the combined lists to obtain
		a common signature poset $P$ and map $\pi\colon Q\to P$,
		\item transport each module's presentation to $P$ by reading off the relevant
		coordinates of the signature.
	\end{enumerate}
	This strategy is implemented by functions such as \texttt{encode\_from\_flanges} in
	\texttt{ZnEncoding.jl} and \texttt{encode\_from\_PL\_fringes} in \texttt{PLPolyhedra.jl}.
	The result is a single poset $P$ that is subordinate to every module in the family.
	
	\paragraph{Theoretical context.}
	The existence of a common encoding for a finite family is built into Miller's
	framework at the level of constant subdivisions, but subtle issues arise if one
	asks for common encodings that capture \emph{all} homomorphisms (full subcategory
	versus tame morphisms).  Waas gives sufficient geometric connectivity conditions
	under which encodable modules form a full abelian subcategory and common encodings
	capture all morphisms \cite{Waa24}.  TAMER-OP's stance is more pragmatic: for a fixed
	finite set of modules and a fixed set of operations (the invariants computed later),
	a common encoding built from the union of the modules' generator regions is the
	appropriate computational object.
	
	\subsection{Complexity and failure modes}
	\label{subsec:encoding-complexity}
	
	Encoding is the stage where combinatorics can blow up.  It is therefore the right
	place to state explicit complexity bounds and to clarify what kinds of inputs are
	expected to be tractable.
	
	\subsubsection*{Complexity of grid-based encodings}
	
	Consider either the lattice or axis-aligned PL backend.
	Let $n$ be the parameter dimension, let $m$ and $r$ be the number of birth and death
	regions in the presentation, and let $C_k$ be the critical coordinate set on axis $k$.
	Define
	\[
	K_k := |C_k|,\qquad
	N_{\mathrm{cell}} := \prod_{k=1}^n (K_k+1),\qquad
	R := |\mathrm{im}(\sigma)|\le N_{\mathrm{cell}}.
	\]
	
	\paragraph{Enumerating signatures.}
	In both backends, the encoding algorithm visits each cell once and performs constant
	work plus updates for constraint events crossing cell boundaries.
	A clean worst-case bound that is easy to compare across variants is:
	\[
	T_{\mathrm{sig}} = O\!\left((m+r)\,n\,N_{\mathrm{cell}}\right),
	\]
	which is the cost of evaluating all membership predicates per cell in a naive scheme.
	The incremental scan improves constants by replacing full reevaluation with sparse
	updates.  Its asymptotic upper bound can be phrased as
	\[
	T_{\mathrm{scan}} = O\!\left(N_{\mathrm{cell}} + E_{\mathrm{flip}}\right),
	\]
	where $E_{\mathrm{flip}}$ is the total number of bit flips processed across the scan.
	In the worst case $E_{\mathrm{flip}} = O(n(m+r)\,N_{\mathrm{cell}})$ because each
	generator bit can flip at most once per axis per fiber of the other coordinates.
	Thus both descriptions agree at the big-O level.
	
	\paragraph{Building the poset order.}
	Given $R$ realized signatures, materializing the order relation by pairwise subset
	tests costs
	\[
	T_{\le} = O\!\left(R^2\cdot \frac{m+r}{w}\right),
	\]
	where $w$ is the machine word size used by the bitset representation (typically $64$).
	If one also computes the Hasse cover relation from the full order matrix (as in
	\texttt{FiniteFringe.\_compute\_cover\_edges\_bitset}), the worst-case bound is
	\[
	T_{\mathrm{cover}} = O\!\left(\frac{R^3}{w}\right),
	\]
	though in practice this is mitigated by bit-parallel operations and by caching cover
	edges only when they are actually requested downstream.
	
	\paragraph{Memory.}
	Storing \texttt{cell\_to\_region} requires $O(N_{\mathrm{cell}})$ memory, which is
	precisely why \texttt{ZnEncodingMap} makes it optional (via \texttt{max\_lookup\_cells}).
	Storing the signature bitvectors requires $O(R(m+r)/w)$ memory.  Materializing the
	full order matrix requires $O(R^2/w)$ memory.
	
	\paragraph{Interpreting the bounds.}
	These formulas highlight the central design tradeoff: for fixed parameter dimension $n$,
	the grid size $N_{\mathrm{cell}}$ is polynomial in the number of critical coordinates
	per axis.  In the lattice backend, each flat or injective contributes at most one
	critical coordinate per constrained axis, so a crude but informative bound is
	$K_k = O(m+r)$, hence
	\[
	N_{\mathrm{cell}} = O\!\left((m+r)^n\right)\quad\text{(for fixed $n$)}.
	\]
	This is the main sense in which finite encodings make computation feasible in the
	common $n=2$ or $n=3$ regimes: the potentially infinite parameter domain is compressed
	to a finite poset whose size grows polynomially in the presentation size when $n$ is
	small.  The same expression also makes the failure mode explicit: for large $n$,
	$(m+r)^n$ is prohibitive, and one must use more problem-specific reductions.
	
	\subsubsection*{Complexity of the polyhedral feasibility backend}
	
	For the polyhedral backend, a different parameterization is appropriate.
	Let $H$ be the total number of halfspaces needed to describe all polyhedral pieces
	in all $U_i$ and $D_j$, and let $B$ measure the branching factor induced by unions
	and complements (formally, the number of piece-selection branches explored).
	
	\paragraph{Feasibility enumeration.}
	The implementation enumerates candidate signatures and branches over polyhedral pieces,
	performing an emptiness check for each branch.  A coarse upper bound is
	\[
	T_{\mathrm{poly}} = O\!\left(2^{m+r}\cdot B \cdot T_{\mathrm{feas}}(n,H)\right),
	\]
	where $T_{\mathrm{feas}}(n,H)$ is the cost of a rational polyhedral feasibility or
	emptiness test in dimension $n$ with $H$ inequalities.  This is exponential in the
	number of generators $m+r$, and this backend is therefore intended for regimes where
	$m+r$ is small but each region has non-axis-aligned geometry.  The implementation
	includes a hard cap \texttt{max\_regions} to prevent runaway enumeration.
	
	\subsubsection*{Typical failure modes}
	
	The most common failure modes are not subtle.
	\begin{itemize}
		\item \textbf{Combinatorial explosion in $N_{\mathrm{cell}}$ or $R$.}
		This is the dominant risk in grid-based encodings and is intrinsic to the problem.
		TAMER-OP mitigates it by offering signature-based collapse (often $R\ll N_{\mathrm{cell}}$)
		and by limiting lookup-table storage.
		\item \textbf{Exponential branching in polyhedral unions.}
		This is intrinsic to complement and union handling in the feasibility backend; the
		mitigation is to use it only when $m+r$ and the number of pieces per region are small,
		and to cap \texttt{max\_regions}.
		\item \textbf{Boundary ambiguity in $\mathbb{R}^n$.}
		The axis-aligned backend avoids this by sampling interior points via \texttt{strict\_eps}.
		The polyhedral backend avoids it by using exact arithmetic for final membership checks.
	\end{itemize}
	
	\subsection{Correctness and diagnostics}
	\label{subsec:encoding-correctness}
	
	The encoding constructors are designed so that correctness follows from a small number
	of verifiable invariants: signatures are computed on cells where they are constant, and
	the poset and transported presentation are defined by those signatures.
	
	\begin{proposition}[Correctness criterion for signature encodings]
		\label{prop:encoding-correctness}
		Suppose $Q$ is partitioned into subsets (cells) $\{C_\alpha\}$ such that for every $i$
		and $j$, the predicates $q\in U_i$ and $q\in D_j$ are constant on each $C_\alpha$.
		Let $P$ be the set of realized signatures on the cells, ordered coordinatewise, and let
		$\pi\colon Q\to P$ send $q$ to the signature of its cell.
		Then $\pi$ is an encoding map subordinate to every indicator module $\kk[U_i]$ and
		$\kk[D_j]$, and any fringe presentation built from these indicators transports to a
		$P$-presentation whose pullback recovers the original one.
	\end{proposition}
	
	\begin{proof}
		By hypothesis, each $U_i$ is a union of cells and each $D_j$ is a union of cells.
		Thus $U_i$ and $D_j$ are unions of fibers of $\pi$, hence they are inverse images of
		subsets of $P$.  The transported subsets are exactly $\widehat{U}_i$ and
		$\widehat{D}_j$ from Subsection~\ref{subsec:constructing-poset}, so
		$\kk[U_i]\cong \pi^*\kk[\widehat{U}_i]$ and similarly for $\kk[D_j]$.
		The statement for a fringe presentation then follows from functoriality: the stalkwise
		matrix description of the map is respected on each fiber because all indicator data are
		constant there.
	\end{proof}
	
	\paragraph{Diagnostics in practice.}
	TAMER-OP exposes and internally relies on a collection of sanity checks that align with
	Proposition~\ref{prop:encoding-correctness}:
	\begin{itemize}
		\item \textbf{Witness-based verification.}  Each realized signature region is stored
		with at least one representative (cell representative in grid backends, polyhedral
		witness point in the feasibility backend).  This makes it straightforward to check that
		the recorded signature agrees with direct membership evaluation.
		\item \textbf{Monotonicity checks.}  Because the intended encoding map is order-preserving
		(Lemma~\ref{lem:signature-monotone}), violations of monotonicity indicate either an
		implementation bug or a boundary-handling issue in $\mathbb{R}^n$; the \texttt{strict\_eps}
		and exact arithmetic mechanisms are designed specifically to prevent such issues.
		\item \textbf{Size and cap reporting.}  Encoding constructors record the numbers
		$N_{\mathrm{cell}}$ and $R$ (or directly $R$ in the feasibility backend) and compare them
		against user-specified caps such as \texttt{max\_regions}.  When a cap is hit, the encoding
		is intentionally aborted rather than silently producing an incomplete poset.
	\end{itemize}
	
	At the thesis level, the main conclusion is that the encoding step is both the
	mathematically essential and the algorithmically delicate part of the TAMER-OP pipeline.
	It is where tameness is converted into finite combinatorics, and it is also where the
	ambient dimension and the geometry of the birth/death regions determine feasibility.
	The next sections treat invariants and derived constructions under the assumption that
	this finite core object $(P,H)$ has been successfully constructed.
	
	
	\section{Invariants}
	
	\subsection{Meta principle: invariants factor through the encoding}
	
	\subsection{Dimension-type invariants}
	
	\subsection{Rank-type invariants}
	
	\subsection{Mobius inversions and signed measures}
	
	\subsection{Slice-based invariants and distances}
	
	\subsection{Vectorizations and kernels}
	
	\subsection{Homological invariants derived from resolutions}
	
	\subsection{Geometry-aware invariants for when $Q = \mathbb{R}^n$}
	
	\subsection{Stability and robustness discussion}
	
	\subsection{Practical API summary}
	
	\section{Homological Algebra and Derived Functors}
	
	\subsection{The Abelian Category of $P$-modules}
	
	\subsection{Projective and Injective Objects}
	
	\subsection{Resolutions}
	
	\subsection{Ext and Tor: definitions and computational model}
	
	\subsection{Multiplicative structures and higher structure}
	
	\subsection{Double complexes and spectral sequences}
	
	\subsection{The ergonomics of derived categories}
	
	\subsection{Relationship to multigraded commutative algebra}
	
	\subsection{What you use all this for in TDA}
	
	\section{Transport and comparison}
	
	\subsection{Change of posets as functorial transport}
	
	\subsection{Comparing modules on different encodings}
	
	\subsection{Comparison via invariants}
	
	\subsection{Interoperability with multipers and RIVET}
	
	\subsection{Case studies}
	
	\section{Statistics}
	
	\subsection{Data model for experiments}
	
	\subsection{Feature extraction recipes}
	
	\subsection{Basic inferential statistics}
	
	\subsection{Learning-oriented workflows}
	
	\subsection{Diagnostic plots}
	
	\subsection{Reproducibility and reporting}
	
	\section{Speed optimizations}
	
	\subsection{Where time is spent}
	
	\subsection{Tradeoffs in exactness versus speed}
	
	\subsection{Sparsity-aware algorithms}
	
	\subsection{Caching and memoization strategies}
	
	\subsection{Parallelism}
	
	\subsection{Geometry backend: controlling combinatorial explosion}
	
	\section{Benchmarks}
	
	\subsection{Philosophy and reproducibility}
	
	\subsection{Microbenchmarks (unit tests)}
	
	\subsection{End-to-end benchmarks}

	\subsection{Comparisons with multipers and RIVET}
	
	\subsection{Discussion of results}
	
	\section{Conclusion}
	aa
	
	% ------------------------------------------------------------
	% References
	% ------------------------------------------------------------

	
	\begin{thebibliography}{99}
		
		\bibitem[Mil20]{Mil20}
		E.~Miller.
		\newblock \emph{Homological algebra of modules over posets}.
		\newblock \href{https://arxiv.org/abs/2008.00063}{arXiv:2008.00063}, 2020.
		
		

		\bibitem[Waa24]{Waa24}
		L.~Waas.
		\newblock \emph{Notes on abelianity of categories of finitely encodable persistence modules}.
		\newblock \href{https://arxiv.org/abs/2407.08666}{arXiv:2407.08666}, 2024.

		\bibitem[Mil23]{Mil23}
		E.~Miller.
		\newblock \emph{Stratifications of real vector spaces from constructible sheaves with conical microsupport}.
		\newblock \emph{Journal of Applied and Computational Topology}, 7(3):473--489, 2023.
		\newblock \href{https://doi.org/10.1007/s41468-023-00112-1}{doi:10.1007/s41468-023-00112-1}.

		\bibitem[KS18]{KS18}
		M.~Kashiwara and P.~Schapira.
		\newblock \emph{Persistent homology and microlocal sheaf theory}.
		\newblock \emph{Journal of Applied and Computational Topology}, 2(1--2):83--113, 2018.
		\newblock \href{https://arxiv.org/abs/1705.00955}{arXiv:1705.00955}.
\bibitem[CB15]{CB15}
		W.~Crawley-Boevey.
		\newblock \emph{Decomposition of pointwise finite-dimensional persistence
			modules}.
		\newblock \emph{Journal of Algebra and Its Applications}, 14(5):1550066, 2015.
		\newblock \href{https://doi.org/10.1142/S0219498815500668}{doi:10.1142/S0219498815500668}.
		
		\bibitem[ZC05]{ZC05}
		A.~Zomorodian and G.~Carlsson.
		\newblock \emph{Computing persistent homology}.
		\newblock \emph{Discrete \& Computational Geometry}, 33:249--274, 2005.
		\newblock \href{https://doi.org/10.1007/s00454-004-1146-y}{doi:10.1007/s00454-004-1146-y}.
		
		\bibitem[CZ09]{CZ09}
		G.~Carlsson and A.~Zomorodian.
		\newblock \emph{The theory of multidimensional persistence}.
		\newblock \emph{Discrete \& Computational Geometry}, 42:71--93, 2009.
		\newblock \href{https://doi.org/10.1007/s00454-009-9176-0}{doi:10.1007/s00454-009-9176-0}.
		
		\bibitem[CEH07]{CEH07}
		D.~Cohen-Steiner, H.~Edelsbrunner, and J.~Harer.
		\newblock \emph{Stability of persistence diagrams}.
		\newblock \emph{Discrete \& Computational Geometry}, 37(1):103--120, 2007.
		\newblock \href{https://doi.org/10.1007/s00454-006-1276-5}{doi:10.1007/s00454-006-1276-5}.
		
		\bibitem[ELZ02]{ELZ02}
		H.~Edelsbrunner, D.~Letscher, and A.~Zomorodian.
		\newblock \emph{Topological persistence and simplification}.
		\newblock \emph{Discrete \& Computational Geometry}, 28:511--533, 2002.
		\newblock \href{https://doi.org/10.1007/s00454-002-2885-2}{doi:10.1007/s00454-002-2885-2}.
		
		\bibitem[Les15]{Les15}
		M.~Lesnick.
		\newblock \emph{The theory of the interleaving distance on multidimensional
			persistence modules}.
		\newblock \emph{Foundations of Computational Mathematics}, 15(3):613--650, 2015.
		\newblock \href{https://doi.org/10.1007/s10208-015-9255-y}{doi:10.1007/s10208-015-9255-y}.
		
		\bibitem[LW15]{LW15}
		M.~Lesnick and M.~Wright.
		\newblock \emph{Interactive visualization of 2-D persistence modules}.
		\newblock \href{https://arxiv.org/abs/1512.00180}{arXiv:1512.00180}, 2015.
		
		\bibitem[RIVET]{RIVET}
		M.~Lesnick and M.~Wright.
		\newblock \emph{RIVET: The rank invariant visualization and exploration tool}
		(software).
		\newblock \url{https://github.com/rivetTDA/rivet}.
		
		\bibitem[LS24]{LS24}
		D.~Loiseaux and H.~Schreiber.
		\newblock \emph{multipers: Multiparameter persistence for machine learning}.
		\newblock \emph{Journal of Open Source Software}, 9(103):6773, 2024.
		\newblock \href{https://doi.org/10.21105/joss.06773}{doi:10.21105/joss.06773}.
		
		\bibitem[CDSGO16]{CDSGO16}
		F.~Chazal, V.~de~Silva, M.~Glisse, and S.~Y. Oudot.
		\newblock \emph{The Structure and Stability of Persistence Modules}.
		\newblock SpringerBriefs in Mathematics. Springer, 2016.
		\newblock \href{https://doi.org/10.1007/978-3-319-42545-0}{doi:10.1007/978-3-319-42545-0}.
		
		\bibitem[Oud15]{Oud15}
		S.~Y. Oudot.
		\newblock \emph{Persistence Theory: From Quiver Representations to Data Analysis}.
		\newblock Mathematical Surveys and Monographs, Vol.~209. American Mathematical
		Society, 2015.
		
		\bibitem[EH10]{EH10}
		H.~Edelsbrunner and J.~L. Harer.
		\newblock \emph{Computational Topology: An Introduction}.
		\newblock American Mathematical Society, 2010.
		
		\bibitem[Mac98]{Mac98}
		S.~Mac~Lane.
		\newblock \emph{Categories for the Working Mathematician}.
		\newblock Graduate Texts in Mathematics, Vol.~5. Springer, 2nd edition, 1998.
		\newblock \href{https://doi.org/10.1007/978-1-4757-4721-8}{doi:10.1007/978-1-4757-4721-8}.
		
		\bibitem[Wei94]{Wei94}
		C.~A. Weibel.
		\newblock \emph{An Introduction to Homological Algebra}.
		\newblock Cambridge Studies in Advanced Mathematics, Vol.~38. Cambridge
		University Press, 1994.
		
		\bibitem[Eis95]{Eis95}
		D.~Eisenbud.
		\newblock \emph{Commutative Algebra with a View Toward Algebraic Geometry}.
		\newblock Graduate Texts in Mathematics, Vol.~150. Springer, 1995.
		\newblock \href{https://doi.org/10.1007/978-1-4612-5350-1}{doi:10.1007/978-1-4612-5350-1}.
		
		\bibitem[MS05]{MS05}
		E.~Miller and B.~Sturmfels.
		\newblock \emph{Combinatorial Commutative Algebra}.
		\newblock Graduate Texts in Mathematics, Vol.~227. Springer, 2005.
		\newblock \href{https://doi.org/10.1007/b138602}{doi:10.1007/b138602}.
		
	\end{thebibliography}
	
	
\end{document}